% Canonical Markdown SHA256: ff1bf841bbcc86d6eb9d289db7fb79e302e20d9d236e287203ef01b8823bda25
% Options for packages loaded elsewhere
\PassOptionsToPackage{unicode}{hyperref}
\PassOptionsToPackage{hyphens}{url}
\PassOptionsToPackage{dvipsnames,svgnames,x11names}{xcolor}
\documentclass[
  english,
  10pt,
  a4paper,
]{article}
\usepackage{xcolor}
\usepackage{amsmath,amssymb}
\setcounter{secnumdepth}{-\maxdimen} % remove section numbering
\usepackage{iftex}
\ifPDFTeX
  \usepackage[T1]{fontenc}
  \usepackage[utf8]{inputenc}
  \usepackage{textcomp} % provide euro and other symbols
\else % if luatex or xetex
  \usepackage{unicode-math} % this also loads fontspec
  \defaultfontfeatures{Scale=MatchLowercase}
  \defaultfontfeatures[\rmfamily]{Ligatures=TeX,Scale=1}
\fi
\usepackage{lmodern}
\ifPDFTeX\else
  % xetex/luatex font selection
\fi
% Use upquote if available, for straight quotes in verbatim environments
\IfFileExists{upquote.sty}{\usepackage{upquote}}{}
\IfFileExists{microtype.sty}{% use microtype if available
  \usepackage[]{microtype}
  \UseMicrotypeSet[protrusion]{basicmath} % disable protrusion for tt fonts
}{}
\makeatletter
\@ifundefined{KOMAClassName}{% if non-KOMA class
  \IfFileExists{parskip.sty}{%
    \usepackage{parskip}
  }{% else
    \setlength{\parindent}{0pt}
    \setlength{\parskip}{6pt plus 2pt minus 1pt}}
}{% if KOMA class
  \KOMAoptions{parskip=half}}
\makeatother
\usepackage{color}
\usepackage{fancyvrb}
\newcommand{\VerbBar}{|}
\newcommand{\VERB}{\Verb[commandchars=\\\{\}]}
\DefineVerbatimEnvironment{Highlighting}{Verbatim}{commandchars=\\\{\}}
% Add ',fontsize=\small' for more characters per line
\newenvironment{Shaded}{}{}
\newcommand{\AlertTok}[1]{\textcolor[rgb]{1.00,0.00,0.00}{\textbf{#1}}}
\newcommand{\AnnotationTok}[1]{\textcolor[rgb]{0.38,0.63,0.69}{\textbf{\textit{#1}}}}
\newcommand{\AttributeTok}[1]{\textcolor[rgb]{0.49,0.56,0.16}{#1}}
\newcommand{\BaseNTok}[1]{\textcolor[rgb]{0.25,0.63,0.44}{#1}}
\newcommand{\BuiltInTok}[1]{\textcolor[rgb]{0.00,0.50,0.00}{#1}}
\newcommand{\CharTok}[1]{\textcolor[rgb]{0.25,0.44,0.63}{#1}}
\newcommand{\CommentTok}[1]{\textcolor[rgb]{0.38,0.63,0.69}{\textit{#1}}}
\newcommand{\CommentVarTok}[1]{\textcolor[rgb]{0.38,0.63,0.69}{\textbf{\textit{#1}}}}
\newcommand{\ConstantTok}[1]{\textcolor[rgb]{0.53,0.00,0.00}{#1}}
\newcommand{\ControlFlowTok}[1]{\textcolor[rgb]{0.00,0.44,0.13}{\textbf{#1}}}
\newcommand{\DataTypeTok}[1]{\textcolor[rgb]{0.56,0.13,0.00}{#1}}
\newcommand{\DecValTok}[1]{\textcolor[rgb]{0.25,0.63,0.44}{#1}}
\newcommand{\DocumentationTok}[1]{\textcolor[rgb]{0.73,0.13,0.13}{\textit{#1}}}
\newcommand{\ErrorTok}[1]{\textcolor[rgb]{1.00,0.00,0.00}{\textbf{#1}}}
\newcommand{\ExtensionTok}[1]{#1}
\newcommand{\FloatTok}[1]{\textcolor[rgb]{0.25,0.63,0.44}{#1}}
\newcommand{\FunctionTok}[1]{\textcolor[rgb]{0.02,0.16,0.49}{#1}}
\newcommand{\ImportTok}[1]{\textcolor[rgb]{0.00,0.50,0.00}{\textbf{#1}}}
\newcommand{\InformationTok}[1]{\textcolor[rgb]{0.38,0.63,0.69}{\textbf{\textit{#1}}}}
\newcommand{\KeywordTok}[1]{\textcolor[rgb]{0.00,0.44,0.13}{\textbf{#1}}}
\newcommand{\NormalTok}[1]{#1}
\newcommand{\OperatorTok}[1]{\textcolor[rgb]{0.40,0.40,0.40}{#1}}
\newcommand{\OtherTok}[1]{\textcolor[rgb]{0.00,0.44,0.13}{#1}}
\newcommand{\PreprocessorTok}[1]{\textcolor[rgb]{0.74,0.48,0.00}{#1}}
\newcommand{\RegionMarkerTok}[1]{#1}
\newcommand{\SpecialCharTok}[1]{\textcolor[rgb]{0.25,0.44,0.63}{#1}}
\newcommand{\SpecialStringTok}[1]{\textcolor[rgb]{0.73,0.40,0.53}{#1}}
\newcommand{\StringTok}[1]{\textcolor[rgb]{0.25,0.44,0.63}{#1}}
\newcommand{\VariableTok}[1]{\textcolor[rgb]{0.10,0.09,0.49}{#1}}
\newcommand{\VerbatimStringTok}[1]{\textcolor[rgb]{0.25,0.44,0.63}{#1}}
\newcommand{\WarningTok}[1]{\textcolor[rgb]{0.38,0.63,0.69}{\textbf{\textit{#1}}}}
\usepackage{longtable,booktabs,array}
\usepackage{calc} % for calculating minipage widths
% Correct order of tables after \paragraph or \subparagraph
\usepackage{etoolbox}
\makeatletter
\patchcmd\longtable{\par}{\if@noskipsec\mbox{}\fi\par}{}{}
\makeatother
% Allow footnotes in longtable head/foot
\IfFileExists{footnotehyper.sty}{\usepackage{footnotehyper}}{\usepackage{footnote}}
\makesavenoteenv{longtable}
\ifLuaTeX
\usepackage[bidi=basic,shorthands=off,]{babel}
\else
\usepackage[bidi=default,shorthands=off,]{babel}
\fi
\ifLuaTeX
  \usepackage{selnolig} % disable illegal ligatures
\fi
\setlength{\emergencystretch}{3em} % prevent overfull lines
\providecommand{\tightlist}{%
  \setlength{\itemsep}{0pt}\setlength{\parskip}{0pt}}
% House preamble for the projection-obstruction paper line.
%
%  Typesetting is not evidence for any claim.
%
% Reviewed 2026-10-02. The `paper-build` skill names a shared
% obstruction-paper-kit preamble; no such file exists in this tree, so this
% deposit carries its own and the archive is self-contained.
%
% xelatex, not pdflatex: the manuscript uses twelve non-ASCII characters
% (section sign, en and em dash, superscript one and two, script ell, capital
% Pi, omega, subscript one, less-or-equal, greater-or-equal, nu) that pdflatex
% cannot set without per-character declarations.

\usepackage{etoolbox}
\usepackage{fontspec}
\usepackage{unicode-math}
\setmainfont{Latin Modern Roman}
% Latin Modern Mono, not DejaVu or Consolas: it ships with every
% TeX distribution, so the archive builds on a machine that has
% only TeX Live or MiKTeX and no extra fonts installed.
\setmonofont{Latin Modern Mono}[Scale=0.85]
\setmathfont{Latin Modern Math}

\usepackage{amsmath,amssymb,amsthm}
\providecommand{\llbracket}{\lBrack}
\providecommand{\rrbracket}{\rBrack}
\usepackage[margin=1in]{geometry}
\usepackage{longtable,booktabs,array}
\usepackage{microtype}
\usepackage{enumitem}
\setlist{nosep,leftmargin=1.4em}

\usepackage[dvipsnames]{xcolor}
\usepackage[colorlinks=true,linkcolor=RoyalBlue,urlcolor=RoyalBlue,
            citecolor=RoyalBlue,breaklinks=true]{hyperref}
% fvextra, not plain fancyvrb: `breaklines` is an fvextra option, and
% fancyvrb alone answers "Package keyval Error: breaklines undefined".
% pandoc defines the Highlighting environment itself, so this preamble sets
% the options globally rather than redefining the environment and racing it.
\usepackage{fvextra}
\usepackage{framed}
\fvset{breaklines=true,breakanywhere=true,fontsize=\small}

% The manuscript tables contain long claim and evidence identifiers. Small type and a
% relaxed tolerance keep them inside the text block instead of overflowing.
\AtBeginEnvironment{longtable}{\small}
\setlength{\tabcolsep}{4pt}
\renewcommand{\arraystretch}{1.15}

% Pandoc emits \tightlist; define it so the run does not fail when the
% template that normally supplies it is not used.
\providecommand{\tightlist}{%
  \setlength{\itemsep}{0pt}\setlength{\parskip}{0pt}}

% Blockquotes carry the paper's refusal boxes and vocabulary notes, so they
% need to read as set-apart rather than merely indented.
\definecolor{quoterule}{gray}{0.75}
\renewenvironment{quote}
  {\begin{leftbar}\small}
  {\end{leftbar}}

% Long file and declaration names are the paper's main source of overfull
% boxes: `po000_l003_double_commutator_estimates.lean` and friends are set in
% typewriter, which by default neither hyphenates nor breaks, so a name wider
% than the remaining line runs into the margin. [htt] lets typewriter text
% hyphenate; the underscore and the other separators below give TeX explicit
% break points inside an identifier, which is what a reader would break at
% anyway.
\usepackage[htt]{hyphenat}
\usepackage{xurl}
\emergencystretch=3em
\sloppy

\usepackage{etoolbox}
\usepackage{titlesec}
\titleformat{\section}{\normalfont\Large\bfseries}{\thesection}{0.6em}{}
\titleformat{\subsection}{\normalfont\large\bfseries}{\thesubsection}{0.6em}{}

\setcounter{secnumdepth}{0}   % the manuscript numbers its own sections
\setcounter{tocdepth}{2}

\usepackage{fancyhdr}
\pagestyle{fancy}
\fancyhf{}
\setlength{\headheight}{14pt}
\fancyhead[L]{\footnotesize PO-001: Covering-Edge Descent}
\fancyhead[R]{\footnotesize PO-001 \textbullet\ v2.0 \textbullet\ 2026-10-02}
\fancyfoot[C]{\thepage}
\renewcommand{\headrulewidth}{0.4pt}
\usepackage{bookmark}
\IfFileExists{xurl.sty}{\usepackage{xurl}}{} % add URL line breaks if available
\urlstyle{same}
\hypersetup{
  pdftitle={From Covering-Edge Descent Complexes to Functorial Presheaves and Coherence Defects},
  pdfauthor={Jeremy H. Carroll},
  pdflang={en},
  colorlinks=true,
  linkcolor={Maroon},
  filecolor={Maroon},
  citecolor={Blue},
  urlcolor={Blue},
  pdfcreator={LaTeX via pandoc}}

\title{From Covering-Edge Descent Complexes to Functorial Presheaves and
Coherence Defects}
\usepackage{etoolbox}
\makeatletter
\providecommand{\subtitle}[1]{% add subtitle to \maketitle
  \apptocmd{\@title}{\par {\large #1 \par}}{}{}
}
\makeatother
\subtitle{Strict Functorial Completion, Coboundary-Curvature Duality,
and the Finite Obstruction Boundary}
\author{Jeremy H. Carroll}
\date{2026-10-02 · Pre-Print (v2.0)}

\begin{document}
\maketitle

\begin{quote}
\textbf{Vocabulary note.} This paper is written inside the CODA
discipline, which fixes a checklist of structural roles (what counts as
a difference, what carries it, what repairs are licensed, what residue
survives, and what the paper authorises --- here, nothing). That
checklist is declared in full in
\hyperref[appendix-a-mechanism-inventory]{Appendix A}. It is programme
vocabulary, not standard mathematical usage, and no result below depends
on reading it first.
\end{quote}

\subsection{Abstract}\label{abstract}

Covering-edge descent data on a finite poset assigns restriction maps
only to adjacent comparable elements. A strict presheaf extension
requires path independence of all saturated-chain composites
(\textbf{Path A}). Alternatively, declaring maps on every comparable
pair gives the composition defect \(K_{c,b,a}=r_{b,a}r_{c,b}-r_{c,a}\)
(\textbf{Path B}). We show \((d_1d_0s)_{c,b,a}=K_{c,b,a}s(c)\),
characterize the flat locus, and compute the degree-3 boundary of the
section-valued curvature. The operator identity \(\delta_2K=0\) holds
without flatness, whereas \(\delta_2\delta_1T=K\circ T-T\circ K\) can be
nonzero; a generic \(H^2\) quotient therefore remains unavailable.

The finite computational extensions (\textbf{Path C}) include exact
rational distance-to-image LP certificates, an additive normal form
under path-independent cover data, bounded-curvature routing, and
diamond commutator traces. The normal form retains mixed terms on longer
triples. A cost-two distinction ledger records a declared charge rather
than an independently derived operator norm. Simplicial face-poset
examples verify the operator cocycle without identifying it with Regge
deficit angles. Bell marginal restrictions remain flat even for a
contextual PR box: contextuality is measured separately by failure of a
nonnegative normalized global probability extension. An exact rational
enclosure is given for the algebraic number \(2(\sqrt{2}-1)\). The new
observed 2-category, protected linear quotient and weighted path bounds
declare their question-preservation premises. Metric hinge holonomy is
derived for rational normal rotations, with an explicit observation and
winding countermodel. Cyclotomic matrix fixtures certify finite
Wilson-loop identities, gauge covariance, and ordered tetrahedral
product closure.

The October 14 supplement constructs the all-degree strict-chain
operator algebra, its graded Leibniz law, and
\(D^2f=K\smile f-f\smile K\), with \(DK=0\). Vertexwise invertible basis
changes conjugate curvature and cannot rectify a non-flat system. A flat
extension preserving the covers exists exactly when all saturated path
products agree, and is then unique. Exact Rust and independent Python
controls implement a bounded rational profile. A supplementary Lean
kernel verifies the graded parity-algebra identities under explicit
differential-algebra assumptions; the paper's pinned Lean package was
not replayed. The historical manifest remains historical evidence.
Higher coherence, physical spacetime, and generic non-flat cohomology
remain outside these results.

\begin{center}\rule{0.5\linewidth}{0.5pt}\end{center}

\subsection{1. Introduction and
Positioning}\label{introduction-and-positioning}

In the analysis of decentralized, distributed, or factorized
approximations on a finite poset \((P,\le)\), local data are compared
across covering edges by contractive linear maps: \[
r_{b,a}:F(b)\to F(a),\qquad (a,b)\in E_{\mathrm{cov}}(P),
\] where \(E_{\mathrm{cov}}(P)=\{(a,b)\in P\times P:a\prec b\}\). This
data suffices to define a degree-0 difference operator: \[
C^0 \xrightarrow{\ d_0\ } C^1,\qquad (d_0s)(a,b)=r_{b,a}s(b)-s(a),
\] measuring local disagreement across adjacent nodes. However,
covering-edge data alone does not define a restriction map
\(r_{c,a}:F(c)\to F(a)\) for non-covering pairs \(a<c\). Without such
maps, there is neither a presheaf functor nor a higher coboundary
complex.

There are two mathematically honest routes forward:

\begin{enumerate}
\def\labelenumi{\arabic{enumi}.}
\tightlist
\item
  \textbf{Path A (Strict Presheaf Completion):} Assume that all
  saturated chains between comparable pairs yield identical composite
  restrictions. This assumption completes the covering-edge data into a
  genuine functor \(P^{\mathrm{op}}\to\mathbf{Ban}\).
\item
  \textbf{Path B (Curvature Cochain):} Declare comparable-pair
  restrictions directly and measure the defect of chain composition via
  a typed degree-2 curvature cochain \(K\).
\end{enumerate}

The conceptual dividing line is the \textbf{Local Green, Global Red
(LGGR)} principle: local covering-edge checks may be satisfied
everywhere (\(d_0s=0\) on covering edges) while higher chain composition
remains obstructed (\(K\ne0\)). Path A eliminates this obstruction by
hypothesis; Path B records and prices it as structural data.

\textbf{Scope.} The algebraic identities apply to finite posets with
linear restriction maps; the Banach-presheaf formulation additionally
requires contractivity. Python fixtures use finite rational matrices and
do not certify arbitrary Banach norms. No claims are made regarding
infinite-dimensional limits, physical spacetime, or quantum gravity.

\begin{center}\rule{0.5\linewidth}{0.5pt}\end{center}

\subsection{2. Path A: Strict Functorial
Presheaves}\label{path-a-strict-functorial-presheaves}

\subsubsection{2.1 Strict Presheaf
Formulation}\label{strict-presheaf-formulation}

Let \((P,\le)\) be a finite partially ordered set.

\begin{quote}
\textbf{{[}DEFINITION{]} Definition 2.1 --- Functorial Banach Presheaf.}
A functorial Banach presheaf on \(P\) is a contravariant functor
\(F:P^{\mathrm{op}}\to\mathbf{Ban}\). It assigns a Banach space \(F(a)\)
to each element \(a\in P\) and a contractive linear map
\(r_{c,a}:F(c)\to F(a)\) to each comparable pair \(a\le c\), satisfying:
1. Identity law: \(r_{a,a}=\mathrm{id}_{F(a)}\) for all \(a\in P\). 2.
Composition law: \(r_{c,a}=r_{b,a}\circ r_{c,b}\) for all
\(a\le b\le c\) in \(P\).
\end{quote}

\begin{quote}
\textbf{{[}PROVEN, CONDITIONAL{]} Proposition 2.1a --- Strict Completion
from Covering Edges.} Let \(P\) be a finite poset equipped with Banach
fibers \(F(a)\) and contractive linear maps \(r_{b,a}:F(b)\to F(a)\) on
all covering edges \(a\prec b\). If for every comparable pair
\(a\le c\), the composition of restriction maps along every saturated
chain from \(a\) to \(c\) is identical, then there exists a unique
functorial Banach presheaf \(\widehat F:P^{\mathrm{op}}\to\mathbf{Ban}\)
extending the covering-edge data.

\emph{Proof.} For any \(a\le c\), choose any saturated chain
\(a=x_0\prec x_1\prec \cdots \prec x_k=c\). Define
\(\widehat r_{c,a}=r_{x_1,x_0}\circ r_{x_2,x_1}\circ \cdots \circ r_{x_k,x_{k-1}}\),
and set \(\widehat r_{a,a}=\mathrm{id}_{F(a)}\). By assumption,
\(\widehat r_{c,a}\) is independent of the chosen chain, making the
assignment well-defined. Contractivity is preserved because finite
compositions of contractive maps are contractive. For any triple
\(a\le b\le c\), concatenating a saturated chain from \(a\) to \(b\)
with one from \(b\) to \(c\) yields a saturated chain from \(a\) to
\(c\). The composite along this chain equals \(\widehat r_{c,a}\) by
definition and equals \(\widehat r_{b,a}\circ\widehat r_{c,b}\) by
associativity. Uniqueness is immediate since covering edges generate all
saturated chains. \(\square\)
\end{quote}

The finite algebraic kernel is mechanized in Lean 4
(\texttt{StrictCompletion.lean}):
\nolinkurl{strict_completion_from_covering_edges},
\nolinkurl{strict_completion_unique_from_chain},
\nolinkurl{diamond_path_independence_of_common_completion},
\nolinkurl{StrictCompletion.chain_law}, and
\nolinkurl{StrictCompletion.diamond_from_chain_law}.

\subsubsection{2.2 Diamonds and Scalar
Witnesses}\label{diamonds-and-scalar-witnesses}

For an elementary diamond \(a\prec b_1\prec c\) and
\(a\prec b_2\prec c\), strict functoriality forces path-independence:

\begin{quote}
\textbf{{[}PROVEN, CONDITIONAL{]} Lemma 2.2 --- Diamond
Path-Independence.} Under Proposition 2.1a, \[
r_{b_1,a}\circ r_{c,b_1} = r_{b_2,a}\circ r_{c,b_2}.
\]

\emph{Proof.} Both sides equal the unique comparable-pair extension
\(\widehat r_{c,a}\). \(\square\)
\end{quote}

\begin{quote}
\textbf{{[}PROVEN, EXAMPLE{]} Scalar Diamond Witness.} Let
\(F(x)=\mathbb{R}\) on a single diamond \(a<b<c\). Let
\(r_{c,b}=\alpha\), \(r_{b,a}=\beta\), and \(r_{c,a}=\mu\). Then the
composition defect on \(v\in\mathbb{R}\) evaluates to \[
K_{c,b,a}(v)=\beta(\alpha v)-\mu v = (\beta\alpha-\mu)v.
\] The defect vanishes if and only if \(\beta\alpha=\mu\). (Mechanized
in Lean 4 as \nolinkurl{Paper001a.scalar_diamond_curvature}.)
\end{quote}

\subsubsection{2.3 Categorical Interface
Unlock}\label{categorical-interface-unlock}

\begin{quote}
\textbf{{[}STANDARD, CONDITIONAL{]} Proposition 2.3 --- Categorical
Interface Unlock (1-categorical only).} Under Proposition 2.1a, the
system satisfies the axioms of a 1-categorical presheaf, unlocking: 1.
The category of elements \(\int_P F\) (Grothendieck construction). 2.
Functorial colimits and limits. 3. Standard descent diagrams over
subposets.

\emph{Scope fence:} This proposition does not assert an
\(\infty\)-categorical enhancement, a model structure, higher homotopy
coherence, or physical dynamics.
\end{quote}

\begin{center}\rule{0.5\linewidth}{0.5pt}\end{center}

\subsection{3. Path B: Curvature Cochains and Higher
Boundaries}\label{path-b-curvature-cochains-and-higher-boundaries}

Path B declares comparable-pair restrictions directly, making no
path-independence assumptions.

\subsubsection{3.1 Composition Curvature}\label{composition-curvature}

\begin{quote}
\textbf{{[}DEFINITION{]} Definition 3.1 --- Composition Curvature.} For
any comparable triple \(a\le b\le c\) in \(P\), the curvature operator
\(K_{c,b,a}:F(c)\to F(a)\) is defined by \[
K_{c,b,a} = r_{b,a}\circ r_{c,b} - r_{c,a}.
\]
\end{quote}

If \(K_{c,b,a}\ne 0\), the system fails the presheaf composition law on
that triple. \textbf{{[}INTERPRETIVE{]}} We interpret \(K\) as a typed
coherence defect.

\emph{Scope fence:} \(K\) is a degree-2 cochain, not a proven cohomology
class in a full \(H^2\) theory.

\subsubsection{\texorpdfstring{3.2 The Differential Identity
\(d_1 d_0 = K\)}{3.2 The Differential Identity d\_1 d\_0 = K}}\label{the-differential-identity-d_1-d_0-k}

Let the cochain spaces over \(P\) be: \[
C^0 = \prod_{x\in P} F(x),\qquad
C^1 = \prod_{a\le b} F(a),\qquad
C^2 = \prod_{a\le b\le c} F(a).
\] The coboundary operators are defined by: \[
(d_0 s)(a,b) = r_{b,a}s(b) - s(a),
\] \[
(d_1 g)(a,b,c) = r_{b,a}(g(b,c)) - g(a,c) + g(a,b).
\]

\begin{quote}
\textbf{{[}PROVEN, CONDITIONAL{]} Theorem 3.2 --- Coboundary-Curvature
Identity.} For every section \(s\in C^0\) and every comparable triple
\(a\le b\le c\): \[
(d_1(d_0 s))(a,b,c) = K_{c,b,a}(s(c)).
\]

\emph{Proof.} Direct expansion yields: \[
\begin{aligned}
(d_1(d_0 s))(a,b,c)
&= r_{b,a}(r_{c,b}s(c) - s(b)) - (r_{c,a}s(c) - s(a)) + (r_{b,a}s(b) - s(a)) \\
&= r_{b,a}(r_{c,b}s(c)) - r_{b,a}s(b) - r_{c,a}s(c) + s(a) + r_{b,a}s(b) - s(a) \\
&= (r_{b,a}\circ r_{c,b} - r_{c,a})(s(c)) \\
&= K_{c,b,a}(s(c)).
\end{aligned}
\] All intermediate section values \(s(a)\) and \(s(b)\) cancel by
linearity. \(\square\)
\end{quote}

Mechanized in Lean 4 as
\nolinkurl{Paper001a.TripleRestriction.d1_d0_eq_curvature}.

\begin{quote}
\textbf{{[}PROVEN, CONDITIONAL{]} Corollary 3.3 --- Flat Locus
Equivalence.} The following are equivalent: 1. \(K_{c,b,a}=0\) for all
\(a\le b\le c\) (the system is flat). 2. \(d_1(d_0 s)=0\) for all
sections \(s\in C^0\). 3. The restriction system satisfies chain
composition on all triples.

\emph{Proof.} (1) \(\implies\) (2) is immediate from Theorem 3.2.
Conversely, if \(d_1(d_0 s)=0\) for all \(s\), let \(v\in F(c)\) be
arbitrary. Choosing a section with \(s(c)=v\) implies
\(K_{c,b,a}(v)=0\), whence \(K_{c,b,a}=0\). The equivalence of (1) and
(3) is Definition 3.1. \(\square\)
\end{quote}

Mechanized in Lean 4 as
\nolinkurl{TripleRestriction.flat_iff_chainComposable} and
\nolinkurl{TripleRestriction.d1_d0_zero_iff_flat}.

\subsubsection{\texorpdfstring{3.3 Degree-3 Boundary for Curvature
(\(d_2 K\))}{3.3 Degree-3 Boundary for Curvature (d\_2 K)}}\label{degree-3-boundary-for-curvature-d_2-k}

For a quadruple \(a\le b\le c\le d\), define the degree-2 coboundary
\(d_2:C^2\to C^3\) by: \[
(d_2 h)(a,b,c,d) = r_{b,a}(h(b,c,d)) - h(a,c,d) + h(a,b,d) - h(a,b,c).
\] Evaluating on the curvature cochain \(h(x,y,z)=K_{z,y,x}(s(z))\):

\begin{quote}
\textbf{{[}PROVEN, CONDITIONAL{]} Theorem 3.4 --- Finite \(d_2K\)
Boundary.} For any quadruple \(a\le b\le c\le d\) and sections
\(s(d)\in F(d)\), \(s(c)\in F(c)\): \[
(d_2 K(s))(a,b,c,d) = K_{c,b,a}(r_{d,c}s(d)) - K_{c,b,a}(s(c)).
\] Consequently, \(d_2 K\) vanishes if: 1. The lower triple is flat
(\(K_{c,b,a}=0\)). 2. The section is edge-compatible along the top link
(\(r_{d,c}s(d)=s(c)\)).

\emph{Proof.} Expand \(d_2\) on the curvature cochain terms: \[
r_{b,a}K_{d,c,b}(s(d)) - K_{d,c,a}(s(d)) + K_{d,b,a}(s(d)) - K_{c,b,a}(s(c)).
\] Substituting definitions: \[
r_{b,a}(r_{c,b}r_{d,c} - r_{d,b})s(d) - (r_{c,a}r_{d,c} - r_{d,a})s(d) + (r_{b,a}r_{d,b} - r_{d,a})s(d) - K_{c,b,a}(s(c)).
\] Regrouping: \[
(r_{b,a}r_{c,b}r_{d,c} - r_{b,a}r_{d,b} - r_{c,a}r_{d,c} + r_{d,a} + r_{b,a}r_{d,b} - r_{d,a})s(d) - K_{c,b,a}(s(c))
= (r_{b,a}r_{c,b} - r_{c,a})(r_{d,c}s(d)) - K_{c,b,a}(s(c)),
\] which is identically \(K_{c,b,a}(r_{d,c}s(d)) - K_{c,b,a}(s(c))\).
\(\square\)
\end{quote}

Mechanized in Lean 4 as
\nolinkurl{QuadRestriction.d2_curvature_eq_lower_curvature_difference},
\nolinkurl{QuadRestriction.d2_curvature_zero_of_lower_flat}, and
\nolinkurl{QuadRestriction.d2_curvature_zero_of_edge_compatible}.

\subsubsection{3.4 Primitive Difference / LGGR
Boundary}\label{primitive-difference-lggr-boundary}

Let data packets carry a local payload and a global curvature defect:
\(p = (x.\mathrm{payload}, x.\mathrm{curvature})\). Consider the
primitive difference relation comparing only payloads: \[
\delta(x, y) = (x.\mathrm{payload}, y.\mathrm{payload}),\qquad \mathrm{ZeroRel}(x,y) \iff x.\mathrm{payload} = y.\mathrm{payload}.
\]

\begin{quote}
\textbf{{[}PROVEN, CONDITIONAL{]} Theorem 3.5 --- Payload-Curvature LGGR
Boundary.} 1. Local payload observations factor through the payload-only
quotient: there exists \(F\) such that \(F([p]) = p.\mathrm{payload}\).
2. If two packets have identical payloads but differing curvatures
(\(k\ne l\)), no curvature observation factors through the payload-only
quotient. 3. If and only if the equivalence relation preserves both
payload and curvature does curvature descend to the quotient.
\end{quote}

Mechanized in Lean 4 as
\nolinkurl{Paper001a.payload_classifier_factors},
\nolinkurl{Paper001a.curvature_not_factor_through_payload_quotient}, and
\nolinkurl{Paper001a.payload_green_curvature_red}.

\subsubsection{3.5 Operator Cocycle and the Curved-Complex
Boundary}\label{operator-cocycle-and-the-curved-complex-boundary}

To test whether \(K\) can define an invariant cohomology class rather
than only a section-dependent defect, introduce operator-valued cochains
\(C^p(P, \mathbf{Hom})\): - \(C^1(P, \mathbf{Hom})\) assigns
\(T_{b,a}: F(b) \to F(a)\) for each comparable pair \(a \le b\). -
\(C^2(P, \mathbf{Hom})\) assigns \(H_{c,b,a}: F(c) \to F(a)\) for each
triple \(a \le b \le c\). - \(C^3(P, \mathbf{Hom})\) assigns
\(\Omega_{d,c,b,a}: F(d) \to F(a)\) for each quadruple
\(a \le b \le c \le d\).

The operator coboundaries are defined by: \[
\begin{aligned}
(\delta_1 T)_{c,b,a} &= r_{b,a} \circ T_{c,b} - T_{c,a} + T_{b,a} \circ r_{c,b}, \\
(\delta_2 H)_{d,c,b,a} &= r_{b,a} \circ H_{d,c,b} - H_{d,c,a} + H_{d,b,a} - H_{c,b,a} \circ r_{d,c}.
\end{aligned}
\]

\begin{quote}
\textbf{{[}PROVEN, CONDITIONAL{]} Theorem 3.6 --- Unconditional Operator
Cocycle Identity.} For any finite poset \(P\), arbitrary fibers
\(F(x)\), and arbitrary restriction maps \(r_{b,a}\), the curvature
cochain \(K_{c,b,a} = r_{b,a} \circ r_{c,b} - r_{c,a}\) satisfies: \[
(\delta_2 K)(a,b,c,d) = r_{b,a} \circ K_{d,c,b} - K_{d,c,a} + K_{d,b,a} - K_{c,b,a} \circ r_{d,c} \equiv 0
\] identically on every quadruple \(a \le b \le c \le d\), with zero
hypothesis on section edge-compatibility or lower flatness.

\emph{Proof.} Direct expansion into compositions yields: \[
\begin{aligned}
(\delta_2 K)(a,b,c,d)
&= r_{b,a}(r_{c,b}r_{d,c} - r_{d,b}) - (r_{c,a}r_{d,c} - r_{d,a}) + (r_{b,a}r_{d,b} - r_{d,a}) - (r_{b,a}r_{c,b} - r_{c,a})r_{d,c} \\
&= r_{b,a}r_{c,b}r_{d,c} - r_{b,a}r_{d,b} - r_{c,a}r_{d,c} + r_{d,a} + r_{b,a}r_{d,b} - r_{d,a} - r_{b,a}r_{c,b}r_{d,c} + r_{c,a}r_{d,c} \\
&= 0. \quad \square
\end{aligned}
\]
\end{quote}

Mechanized in Lean 4 as
\nolinkurl{Paper001a.QuadRestrictionOp.delta2_operator_curvature_identically_zero}
and verified in Python via \nolinkurl{code/cohomology_quotient.py}.

\begin{quote}
\textbf{{[}PROVEN, EXAMPLE{]} Theorem 3.7 --- Generic \(H^2\) Quotient
Obstruction.} The maps above do not form a cochain complex on an
arbitrary non-flat restriction system. In general, \[
(\delta_2\delta_1T)_{d,c,b,a}
=K_{c,b,a}\circ T_{d,c}-T_{b,a}\circ K_{d,c,b},
\] which need not vanish. For scalar fibers with \(r_{b,a}=2\),
\(r_{c,b}=3\), \(r_{c,a}=0\), \(T_{d,c}=1\), and the other relevant
\(T\) components zero, the residual is \(6\). Therefore
\(\mathrm{im}(\delta_1)\) is not generally contained in
\(\ker(\delta_2)\), so \(\ker(\delta_2)/\mathrm{im}(\delta_1)\) is not a
defined generic quotient for Path B.
\end{quote}

The Python \nolinkurl{OperatorCochainImageDiagnostic} performs exact
membership tests in \(\mathrm{im}(\delta_1)\) only. On a separately
established flat complex, the usual quotient construction may be
considered; no generic gauge-equivalence or rectification theorem is
claimed here.

\begin{center}\rule{0.5\linewidth}{0.5pt}\end{center}

\subsection{4. Finite Extensions, Model Consequences, and Evidence
Boundaries}\label{finite-extensions-model-consequences-and-evidence-boundaries}

Sections 4.1--4.11 retain their existing numbers for traceability. They
distinguish algebraic identities, conditional constructions, declared
model conventions, and finite examples. Passing an example or mutation
test is evidence for that fixture; it is not a proof of every
surrounding interpretation.

\subsubsection{4.1 Curved Differential Derivation Identity (Theorem
4.1)}\label{curved-differential-derivation-identity-theorem-4.1}

In non-flat restriction systems where \(K \ne 0\), the coboundary
operators \(\delta\) do not form a standard nilquadratic cochain complex
(\(\delta^2 \ne 0\)). We show that the defect is governed by the inner
commutator with the curvature cochain.

\begin{quote}
\textbf{{[}PROVEN, CONDITIONAL{]} Theorem 4.1 --- Curved Differential
Derivation Identity.} Let \((P, \le)\) be a finite poset equipped with
Banach fibers \(F(x)\) and declared comparable-pair restriction maps
\(r_{b,a}\). On the operator cochain space \(C^1(P, \mathbf{Hom})\), the
composition \(\delta_2 \delta_1\) satisfies: \[
(\delta_2 \delta_1 T)_{d,c,b,a} = [K, T]_{d,c,b,a} = K_{c,b,a} \circ T_{d,c} - T_{b,a} \circ K_{d,c,b}
\] for all quadruples \(a \le b \le c \le d\), while the curvature
itself satisfies \(\delta_2(K) \equiv 0\). The all-degree strict-chain
construction below supplies the additional product and differential
laws; it does not construct compositor 2-cells.

\emph{Proof.} Expanding \((\delta_2(\delta_1 T))_{d,c,b,a}\) using the
definition of \(\delta_2\) on \(H = \delta_1 T\): \[
(\delta_2 H)_{d,c,b,a} = r_{b,a} \circ (\delta_1 T)_{d,c,b} - (\delta_1 T)_{d,c,a} + (\delta_1 T)_{d,b,a} - (\delta_1 T)_{c,b,a} \circ r_{d,c}.
\] Substituting
\((\delta_1 T)_{z,y,x} = r_{y,x} \circ T_{z,y} - T_{z,x} + T_{y,x} \circ r_{z,y}\)
and collecting terms with respect to \(T\), the cross terms involving
\(r_{y,x} \circ r_{z,y}\) recombine with the direct terms \(r_{z,x}\) to
form \(K_{z,y,x} = r_{y,x} \circ r_{z,y} - r_{z,x}\). All purely linear
terms cancel identically, leaving
\(K_{c,b,a} \circ T_{d,c} - T_{b,a} \circ K_{d,c,b}\). The cocycle
condition \(\delta_2(K) = 0\) follows from Theorem 3.6. \(\square\)
\end{quote}

\emph{Computational verification:} Tested in
\texttt{code/curved\_bicomplex.py} and
\nolinkurl{code/tests/test_curved_bicomplex.py::test_nonflat_diamond_matrix_2x2_fibers}.

\textbf{All-degree construction (PO001-ALL-DEGREE-CURVED).} Fix modules
\(V_x\) over a commutative unital ring and use linear maps, or Banach
fibers and bounded linear maps. Declare the strict-chain carrier \[
C^n=\prod_{x_0<\cdots<x_n}\operatorname{Hom}(V_{x_n},V_{x_0}),\qquad n\ge0.
\] Here degree zero consists of vertex endomorphisms, distinct from the
section-valued \(C^0\) in §3.2. Self-restrictions are identities when
recovering the full presheaf interface; the executable strict-chain
profile has no separate self-edge inputs. Its product, internal-face
differential, and connection are \[
\begin{aligned}
(f\smile g)_{x_0\ldots x_{p+q}}
 &=f_{x_0\ldots x_p}\circ g_{x_p\ldots x_{p+q}},\\
(bf)_{x_0\ldots x_{n+1}}
 &=\sum_{i=1}^{n}(-1)^i f_{x_0\ldots\widehat{x_i}\ldots x_{n+1}},\\
A_{x,y}&=r_{y,x},\qquad
Df=bf+A\smile f-(-1)^n f\smile A,\\
K&=bA+A\smile A.
\end{aligned}
\] The empty internal sum in degree zero is zero. Only internal
deletions occur in \(b\), so endpoint types agree; deleting endpoints
instead would require extra transports. The unit is the degree-zero
family of identity maps. Degrees beyond the maximum strict-chain length
minus one are zero spaces.

\begin{quote}
\textbf{{[}PROVEN, CONDITIONAL{]} Theorem 4.1a --- All-Degree Curved
Operator Algebra.} This is a unital associative graded algebra with \[
b^2=0,\quad b(f\smile g)=bf\smile g+(-1)^p f\smile bg,
\] \[
D(f\smile g)=Df\smile g+(-1)^p f\smile Dg,\qquad
D^2f=K\smile f-f\smile K,\qquad DK=0.
\] Thus \((C^\bullet,b)\) is a differential graded algebra;
\((C^\bullet,D,K)\) is its curved counterpart. It is generally
noncommutative.

\emph{Proof.} Associativity and the unit follow componentwise from
composition. For \(b^2\), every pair of internal deletions \(i<j\)
occurs twice, with signs \((-1)^{i+j-1}\) and \((-1)^{i+j}\), and
cancels. In \(b(f\smile g)\) split deleted vertices at \(i\le p\) and
\(i>p\). The first group is \(bf\smile g\), sharing vertex \(x_{p+1}\);
the second is \((-1)^p f\smile bg\), sharing \(x_p\). This establishes
the Leibniz law also when one degree is zero. Expanding the graded
commutator \([A,f]=A\smile f-(-1)^n f\smile A\) gives its Leibniz law,
and hence that of \(D\). Apply both laws once more: the terms involving
\(bf\) cancel and \(b^2f=0\), leaving \[
D^2f=(bA)\smile f-f\smile(bA)
      +(A\smile A)\smile f-f\smile(A\smile A).
\] Finally \(bK=(bA)\smile A-A\smile bA\), while
\(A\smile K-K\smile A=A\smile bA-(bA)\smile A\). Their sum is \(DK=0\).
All sums are finite; no norm or convergence assumption is used. On
strict triples and quadruples, \(D\) recovers \(\delta_1\) and
\(\delta_2\). \(\square\)
\end{quote}

The same interior-deletion formulas can be placed on weak chains, but
this is a different cochain carrier. Extending a strict-chain cochain by
zero on repeated vertices does not intertwine \(b\): on \((x,x,y)\),
\(bA=-A_{x,y}\) can be nonzero. The implementation and the flatness
detector below use strict chains, without a claim about normalized
weak-chain subcomplexes.

\begin{quote}
\textbf{{[}PROVEN, CONDITIONAL{]} Corollary 4.1b --- Full Operator
Detector of Flatness.} For this carrier, \(K=0\) if and only if \(D^2\)
vanishes on every degree-zero operator cochain; equivalently, \(D^2=0\)
in every degree.

\emph{Proof.} Flatness implies nilpotence by Theorem 4.1a. Conversely,
for each strict triple \(x<y<z\) take the degree-zero vertex selector
\(e_z\), identity at \(z\) and zero elsewhere. Then
\((K\smile e_z-e_z\smile K)_{x,y,z}=K_{x,y,z}\). If every such
\(D^2e_z\) vanishes, every curvature component vanishes. \(\square\)
\end{quote}

The full operator carrier is essential: testing only the global identity
cochain always gives \([K,1]=0\) and does not detect curvature. Over
\(\mathbb Q\), the graded commutator
\([f,g]=f\smile g-(-1)^{pq}g\smile f\) also gives a graded Lie algebra:
expanding the six associative triple products establishes graded Jacobi,
and the Leibniz law makes \(b\) a Lie derivation. The curvature equation
is \(K=bA+\tfrac12[A,A]\), so flatness is the Maurer--Cartan equation in
this declared algebra. This is an algebraic Lie bridge, with no physical
symmetry group inferred.

\emph{Evidence scope.} The written theorem applies to the declared
all-degree carrier on any finite poset.
\nolinkurl{code/operator_incidence_dga.py} is an independent rational
reference, with \nolinkurl{code/tests/test_operator_incidence_dga.py}
controls. The reusable Rust owner is
\nolinkurl{omega_homology::restriction_dga}; its executable profile has
at most eight vertices and fiber dimensions at most four, complete
component coverage, and exact rational arithmetic. Nonzero \(D^2\)
controls occur in degrees zero through four on an eight-vertex chain.
\nolinkurl{lean4/AllDegreeCurvature.lean} retains the earlier
assumed-law parity kernel. \nolinkurl{lean4/ConcreteIncidence.lean}
constructs the recursive signed internal deletion, overlapping cup and
covariant derivative over an arbitrary common noncommutative coefficient
ring. It constructs strict-poset chain components and establishes their
intrinsic positive-degree nilpotence, graded Leibniz and curvature
square from those definitions; degree-zero endpoint actions have a
separate concrete square theorem. The varying endpoint-Hom construction
is now faithfully embedded by
\nolinkurl{lean4/HeterogeneousIncidence.lean}, including intrinsic
strict-poset operations and the vertex selector; graded Jacobi remains a
written proof with native controls, outside this Lean replay. The
supplementary reproduction profile and printed axiom dependencies are
recorded in \href{docs/derivations_2026_10_14.md}{the derivation
review}.

\subsubsection{4.1c Faithful heterogeneous endpoint-Hom
transport}\label{c-faithful-heterogeneous-endpoint-hom-transport}

For modules \(V_x\) over a commutative ring, set \(E=\prod_xV_x\) and
let \(i_x:V_x\to E\) and \(p_x:E\to V_x\) be the coordinate injection
and projection. For finite \(P\), \(E\) is also the finite direct sum.
Define \[
L_{xz}(f)=i_x f p_z\in\operatorname{End}(E).
\] The identities \(p_xi_x=I\) and \(p_xi_z=0\) for \(x\ne z\) give \[
p_xL_{xz}(f)i_z=f,\qquad
L_{xy}(f)L_{yz}(g)=L_{xz}(fg).
\] Thus every corner map is injective, addition is preserved, and
internal face deletion commutes with \(L\) because its endpoints are
fixed. Each cup cut has a common pivot, so \(L(f\smile g)=Lf\smile Lg\).
Consequently the already constructed common-ring deletion and covariant
identities transfer faithfully to the actual varying Hom fibers. Strict
subchains remain strict; embedding also commutes with zero extension and
restriction. This establishes intrinsic strict-poset \(b^2=0\) and
\(D^2=[K,-]\) without adding differential laws as assumptions.

Let \(e_x=L_{xx}(I)\). These are orthogonal idempotents. On a strict
triple \(x<y<z\), the vertex selector supported at \(z\) gives
\(L_{xz}(K)e_z-0L_{xz}(K)=L_{xz}(K)\), so a vanishing selected component
forces \(K=0\) by extraction. Degree-zero endpoint square is supplied by
the concrete kernel; the selector identity is proved in the
heterogeneous kernel.

\emph{Scope.} \nolinkurl{HeterogeneousIncidence.lean} establishes these
injection, composition, transport, strict-poset square and detector
identities for arbitrary modules with the declared types. The finite
Rust corner owner is \nolinkurl{omega_homology::block_embedding},
retaining the original carrier identity. The six-term graded Jacobi
expansion in §4.1 is a written mathematical proof and an exact native
control; it is not a new Lean theorem. Neither the embedding nor Jacobi
declares a physically available Lie action.

\subsubsection{4.2 Four-Term Operator Comparison and the
Higher-Coherence Boundary (Theorem
4.2)}\label{four-term-operator-comparison-and-the-higher-coherence-boundary-theorem-4.2}

\begin{quote}
\textbf{{[}PROVEN, CONDITIONAL{]} Theorem 4.2 --- Four-Term Operator
Comparison.} For declared linear restrictions on \(a\le b\le c\le d\),
\[
r_{b,a}K_{d,c,b}+K_{d,b,a}
=K_{d,c,a}+K_{c,b,a}r_{d,c}.
\] \emph{Proof.} Substitute \(K_{z,y,x}=r_{y,x}r_{z,y}-r_{z,x}\). Both
sides reduce, by associativity, to \(r_{b,a}r_{c,b}r_{d,c}-r_{d,a}\).
This is Theorem 3.6 rearranged. \(\square\)
\end{quote}

All terms are maps \(F(d)\to F(a)\). The identity compares two ways of
accumulating operator defects. It does not supply categories as fibers,
natural transformations as compositor 2-cells, units, or higher
coherence. In particular, \(K_{c,b,a}:F(c)\to F(a)\) need not be an
endomorphism, so an expression \(I+K_{c,b,a}\) is generally not typed.
The following realization adds a specific target 2-category and
admission premises; the four-term identity alone supplies neither.

\emph{Computational verification:}
\nolinkurl{code/tests/test_cohomology_quotient.py::test_operator_cocycle_identity_linear_chain}.

\subsubsection{4.2a A typed observed 2-category and normal lax
functor}\label{a-a-typed-observed-2-category-and-normal-lax-functor}

Declare an object \((V,Q)\) with a vector space and a finite family of
linear questions \(Q:V\to A_Q\). Put \(N_Q=\ker Q\). A 1-cell
\(f:(V,Q)\to(W,R)\) is a linear map satisfying \(fN_Q\subseteq N_R\).
There is exactly one 2-cell \(f\Rightarrow g\) between parallel 1-cells
when \(R(g-f)=0\), and no 2-cell otherwise. Its retained label is
\(\alpha=g-f\). Vertical composition adds differences; the identity
label is zero. For \(\beta:g\Rightarrow g'\) and
\(\alpha:f\Rightarrow f'\), \[
\beta*\alpha=g'f'-gf=g\alpha+\beta f+\beta\alpha.
\] The mixed term is essential. The label is invisible because
\(\alpha\) has image in \(N_R\), \(g'\) preserves that kernel, and
\(\beta f\) is invisible at the final object. This establishes
horizontal closure; omitting the 1-cell kernel condition can expose a
formerly invisible distinction. Matrix associativity gives associative
1-composition; addition and uniqueness give vertical associativity,
units, horizontal associativity and interchange. Hence these objects,
arrows and cells form a specified strict, locally thin 2-category. The
cells are invertible comparisons of declared answers, with no inference
of physically reversible dynamics or a natural transformation between
unspecified fiber categories.

View \(P^{op}\) as a locally discrete 2-category. Give every vertex its
declared \(Q_x\), set \(r_{xx}=I\), require every \(r_{xy}\) to preserve
question kernels, and require \(Q_xK_{xyz}=0\) on every strict triple.
Then the assignments have compositors \[
\phi_{xyz}:r_{xy}r_{yz}\Rightarrow r_{xz},\qquad
\phi_{xyz}=r_{xz}-r_{xy}r_{yz}=-K_{xyz}.
\] Self-edge unit cells have label zero. On a quadruple, both composites
are parallel and have the same label: \[
(r_{xz}-r_{xy}r_{yz})r_{zw}+(r_{xw}-r_{xz}r_{zw})
=r_{xy}(r_{yw}-r_{yz}r_{zw})+(r_{xw}-r_{xy}r_{yw})
=r_{xw}-r_{xy}r_{yz}r_{zw}.
\] Every five-way pentagon commutes because its paths are admitted
parallel cells in a locally thin hom-category. This is a normal lax
functor into the specified target. Since its comparison cells are
invertible, it is also a pseudofunctor into that target. With full
separating questions, a cell exists only for equal operators and these
premises force \(K=0\). With zero questions, all parallel maps compare:
that codiscrete envelope loses all protected distinctions. Intermediate
questions can retain nonzero invisible curvature, as the exact controls
show. A rectangular \(K\) is a difference label with endpoints, never an
untyped \(I+K\). Arbitrary visible curvature is refused.

\emph{Scope.} The construction and complete coherence proof are written
above. \nolinkurl{omega_homology::observed_2category} verifies the
actual carrier identities, 1-cells and every compositor.
\nolinkurl{code/observed_transport.py} independently checks units,
interchange, pentagon and refusals. \texttt{TransportIdentities.lean}
replays the ring comparison kernels; it does not formalize the entire
typed 2-category. This realizes one observer-sensitive comparison
semantics, not all possible higher-cell extensions or native Genesis
dynamics.

\subsubsection{4.3 Exact Primal-Dual Distance-to-Image Certification
(Theorem
4.3)}\label{exact-primal-dual-distance-to-image-certification-theorem-4.3}

Choose bases in the finite operator-cochain spaces and write \(D\) for
the rational matrix of \(\delta_1\). For a rational cochain \(Z\)
(including \(Z=K\)), define \[
D_1(Z)=\min_t\|Z-Dt\|_1,
\] where the norm is the sum of absolute coordinate entries in these
declared bases. This is a distance to a linear image, not a generic
\(H^2\) class or a theorem about nonlinear gauge transformations.

\begin{quote}
\textbf{{[}PROVEN, CONDITIONAL{]} Theorem 4.3 --- Exact Primal-Dual
Certification.} The primal and dual programs \[
\min_{t,u}\mathbf{1}^Tu\quad\text{subject to}\quad -u\le Z-Dt\le u,
\] \[
\max_\lambda\lambda^TZ\quad\text{subject to}\quad D^T\lambda=0,\quad\|\lambda\|_\infty\le1
\] attain the same rational value. An independently checked feasible
primal-dual pair with equal objectives certifies this optimum.

\emph{Proof.} The primal is feasible at \(t=0,u=|Z|\) and bounded below
by zero. The dual is feasible at zero. Finite linear-programming duality
gives attainment and equal objectives, and rational data admit rational
optimal solutions. Weak duality makes a feasible equal-objective pair a
directly checkable certificate. \(\square\)
\end{quote}

\emph{Computational verification:} \nolinkurl{code/exact_coda_solver.py}
(\texttt{ExactCodaL1Solver}) and
\nolinkurl{code/tests/test_coda_dual_witness.py::test_obstructed_cochain_positive_phi1_zero_gap}.
These tests include arbitrary cochains; a positive certificate for one
cochain does not assert positivity for every curvature cochain. The
certificate establishes the declared finite LP, with no cryptographic
authenticity claim.

\subsubsection{4.4 Conditional Additive Normal Form (Theorem
4.4)}\label{conditional-additive-normal-form-theorem-4.4}

\begin{quote}
\textbf{{[}PROVEN, CONDITIONAL{]} Theorem 4.4 --- Cover-Preserving
Additive Normal Form.} Suppose \(r\) is declared on all comparable pairs
of a finite poset, with identity maps on self-pairs, and its covering
maps satisfy the path-independence hypothesis of Proposition 2.1a. Let
\(\widehat r\) be their unique strict completion and set
\(A_{c,a}=r_{c,a}-\widehat r_{c,a}\). Then \(r=\widehat r+A\), \(A=0\)
on covers and identities, and \[
\begin{aligned}
K_{c,b,a}={}&\widehat r_{b,a}A_{c,b}+A_{b,a}\widehat r_{c,b}\\
&+A_{b,a}A_{c,b}-A_{c,a}.
\end{aligned}
\] If both \(a\prec b\) and \(b\prec c\) are covering pairs, this
reduces to \(K_{c,b,a}=-A_{c,a}\).

\emph{Proof.} Proposition 2.1a constructs \(\widehat r\) under the
stated hypothesis. Expand
\((\widehat r_{b,a}+A_{b,a})(\widehat r_{c,b}+A_{c,b})-(\widehat r_{c,a}+A_{c,a})\)
and cancel the flat-background terms. The cover-pair specialization sets
both inner corrections to zero. \(\square\)
\end{quote}

This is an additive decomposition; no orthogonality or direct-sum
structure is asserted. The path-independence hypothesis is necessary for
a flat background preserving every cover. On the contractive scalar
diamond \(a\prec x,y\prec d\), the choices \(r_{x,a}=r_{d,x}=1/2\),
\(r_{y,a}=1/3\), and \(r_{d,y}=1/2\) give path products \(1/4\) and
\(1/6\). No such background exists. Selecting a spanning forest cannot
repair the omitted path.

On the longer chain \(a\prec b\prec c\prec d\) with covers \(1/2\),
direct maps \(r_{c,a}=r_{d,b}=1/3\), and \(r_{d,a}=1/10\), the triple
\((a,b,d)\) has \(K_{d,b,a}=1/15\), whereas \(-A_{d,a}=1/40\). The mixed
terms are essential.

\emph{Computational verification:}
\nolinkurl{code/curved_bicomplex.py::decompose_normal_form} refuses
noncommuting cover data. The exact positive, refusal, and longer-triple
tests are in \nolinkurl{code/tests/test_review_boundaries.py}.

\textbf{Declared gauge and rectification boundary
(PO001-GAUGE-RECTIFICATION).} Use the vertexwise change-of-basis group
\(G=\prod_x\operatorname{GL}(V_x)\) of linear automorphisms (bounded
with bounded inverses for Banach fibers), with declared inverses
\(h_x=g_x^{-1}\), acting by \[
r^g_{y,x}=g_xr_{y,x}h_y,\qquad
f^g_{x_0\ldots x_n}=g_{x_0}f_{x_0\ldots x_n}h_{x_n}.
\] These are algebraic basis changes. General linear transformations
need not preserve fixed Banach norms; if contractivity in those norms is
required, restrict to admitted isometries or transport the norms. No
operational availability or unitary interpretation is assumed.

\begin{quote}
\textbf{{[}PROVEN, CONDITIONAL{]} Theorem 4.4a --- Gauge No-Go and
Unique Flat Cover Extension.} The declared action preserves the product
and intertwines \(b\) and \(D\): \[
(f\smile u)^g=f^g\smile u^g,\qquad
(bf)^g=b(f^g),\qquad
(D_Af)^g=D_{A^g}(f^g),\qquad
K(r^g)_{x,y,z}=g_xK(r)_{x,y,z}h_z.
\] Consequently a gauge orbit contains a flat representative if and only
if the original restriction system is flat. Separately, prescribed cover
maps have a cover-preserving flat extension if and only if every
saturated-chain product with the same endpoints agrees; this extension
is unique, with identities on self-pairs.

\emph{Proof.} In a cup product the adjacent factors \(h_{x_p}g_{x_p}\)
cancel. Internal deletions leave both endpoints unchanged, so \(b\)
commutes with the action. These two facts yield the \(D\) and \(K\)
formulas. Multiplying a zero transformed curvature on the left by
\(h_x\) and on the right by \(g_z\) recovers the original component;
hence zero curvature is invariant. For the extension criterion,
necessity follows by applying functoriality to each saturated chain. For
sufficiency, define each comparable-pair map to be its common
saturated-chain product and identities to be identity maps.
Concatenation of saturated chains establishes composition; generation by
covers yields uniqueness, as in Proposition 2.1a. This compares all
saturated paths, not only selected diamonds. \(\square\)
\end{quote}

The scalar diamond above is an exact obstruction to cover-preserving
flat extension: \(1/4\ne1/6\). Conversely, commuting covers can admit
\(\widehat r\) while the independently declared long-edge maps of \(r\)
make \(K(r)\ne0\). Replacing those long-edge maps by \(\widehat r\)
removes curvature by changing the restriction system. It is not an
invertible gauge transformation of that curved system. The expansion in
Theorem 4.4 retains both mixed terms.

Nor does a norm optimum select a canonical representative. For
\(w=(1,0)\) and repair line \(\operatorname{span}(1,1)\), \[
\|w-t(1,1)\|_1=|1-t|+|t|\ge1,
\] with equality for every \(0\le t\le1\). A distance-to-image
certificate therefore supplies neither a unique repair nor gauge
rectification. The curved algebra also leaves Theorem 3.7 intact: an
ordinary \(H^2\) quotient requires
\(\operatorname{im}\delta_1\subseteq\ker\delta_2\) in the actual chosen
complex. Other gauge actions, noninvertible coarse-grainings, and
physical rectification mechanisms require their own definitions and are
not classified here.

\emph{Evidence scope.}
\nolinkurl{code/tests/test_gauge_rectification.py} independently checks
the declared action, flat extension, path obstruction, mixed terms,
malformed inverses, and nonunique norm minimizers.
\nolinkurl{omega_homology::restriction_dga} checks inverse identities
and returns disagreeing paths rather than admitting incomplete cover
data. The supplementary Lean kernel shows only that invertible endpoint
factors preserve zero; it does not mechanize arbitrary-poset
saturated-path enumeration or physical gauge availability.

\subsubsection{4.5 Prime Triad Descent Soundness and Bounded-Curvature
Routing (Theorem
4.5)}\label{prime-triad-descent-soundness-and-bounded-curvature-routing-theorem-4.5}

Finite operational packet routing across curved posets is governed by
the Omega Engine Algorithm Prime Triad: - \textbf{Algorithm Prime -1:}
Grounding of primitives and carrier honesty. - \textbf{Algorithm Prime
000:} Local covering-edge declaration order (\(a \prec b\)). -
\textbf{Algorithm Prime 001:} Adapter descent soundness verification. -
\textbf{Algorithm Prime 002:} Continue-until-unauthorized control law
with curvature budget \(\kappa \le \kappa_{\max}\).

\begin{quote}
\textbf{{[}PROVEN, CONDITIONAL{]} Theorem 4.5 --- Prime Triad
Bounded-Curvature Routing Law.} Let a packet \(\pi\) traverse a path
\(\gamma = (x_0 \prec x_1 \prec \cdots \prec x_k)\) in \(P\). Under
Algorithm Prime 002, the packet is authorized to proceed if and only if:
1. Every step is an elementary covering edge (Prime 000). 2. The adapter
map satisfies descent soundness
\(\|r_{x_{i}, x_{i-1}} v - v_{i-1}\| \le \epsilon\) (Prime 001). 3. The
accumulated path curvature does not exceed the declared threshold: \[
   \kappa(\gamma) = \sum_{i=1}^{k-1} \|K_{x_{i+1}, x_i, x_{i-1}}\|_1 \le \kappa_{\max}.
   \] If \(\kappa(\gamma) > \kappa_{\max}\), Prime 002 halts
immediately, revokes transmission authority, and logs a hashed tombstone
in the cemetery DAG.

\emph{Proof.} Immediate from the operational definition of the Prime 002
state machine. The hash-chain checker detects inconsistencies in its
stored receipt fields. It does not authenticate independently supplied
measurements or prevent an actor from recomputing hashes. \(\square\)
\end{quote}

\emph{Computational verification:} Implemented in
\nolinkurl{code/prime_triad_router.py} and tested in
\nolinkurl{code/tests/test_prime_triad_router.py}.

\subsubsection{4.6 DDC Operational Commutator on Poset Lattices (Theorem
4.6)}\label{ddc-operational-commutator-on-poset-lattices-theorem-4.6}

In Discrete Deformation Calculus (DDC), discrete state transitions are
encoded using the 8-letter canonical rune alphabet: \[
\mathcal{A}_{\mathrm{DDC}} = \{E, R, C, I, \Lambda, D, M, T\},
\] representing Expansion (\(E\)), Reduction (\(R\)), Contraction
(\(C\)), Insertion (\(I\)), Lift (\(\Lambda\)), Descent (\(D\)),
Mutation (\(M\)), and Terminal (\(T\)).

\begin{quote}
\textbf{{[}PROVEN, CONDITIONAL{]} Theorem 4.6 --- DDC Operational
Commutator Theorem.} On an elementary diamond lattice
\(a \prec b_1, b_2 \prec c\), let \(\Lambda_{b_1, a}\) and
\(D_{a, b_1}\) denote the lift and descent operations along the left
branch, and \(\Lambda_{b_2, a}, D_{a, b_2}\) along the right branch. The
operational commutator: \[
\llbracket \Lambda, D \rrbracket = D_{a, b_1} \circ \Lambda_{c, b_1} - D_{a, b_2} \circ \Lambda_{c, b_2}
\] evaluates identically to the diamond holonomy defect
\(K_{c, b_1, a} - K_{c, b_2, a}\). The commutator vanishes if and only
if the diamond is flat.

\emph{Proof.} Expanding the definitions of lift and descent as
restriction and extension maps on poset sections:
\(D_{a, b} \circ \Lambda_{c, b} = r_{b, a} \circ r_{c, b}\). Subtraction
across the two branches gives
\((r_{b_1, a} \circ r_{c, b_1} - r_{c, a}) - (r_{b_2, a} \circ r_{c, b_2} - r_{c, a}) = K_{c, b_1, a} - K_{c, b_2, a}\).
\(\square\)
\end{quote}

\emph{Computational verification:} Implemented in
\texttt{code/ddc\_poset\_trace.py} and tested in
\nolinkurl{code/tests/test_ddc_and_downmap_mutations.py::test_ddc_diamond_commutator_flat_vs_curved}.

\subsubsection{4.7 Distinction Collapse and the Declared Cost-Two Ledger
(Model Consequence
4.7)}\label{distinction-collapse-and-the-declared-cost-two-ledger-model-consequence-4.7}

The finite down-map model declares a set of protected distinction
records and charges \(\mu_1=2\) for each active record whose endpoints
receive the same image.

\begin{quote}
\textbf{{[}MODEL CONSEQUENCE{]} Statement 4.7 --- Declared Ledger
Charge.} If \(m\) active protected records collapse under the supplied
down-map, its ledger entry has \[
\mathrm{loss\_count}=m,\qquad \mathrm{loss\_cost\_mu1}=2m.
\] \emph{Proof.} The implementation counts these records and assigns
their count times the declared cost. \(\square\)
\end{quote}

This convention does not derive \(\|K\|_1\ge2m\) for an independent
operator curvature. The legacy field \texttt{induced\_curvature\_norm}
stores the same declared charge; it is not a second measurement.
\nolinkurl{verify_antimorphic_curvature_lower_bound} checks a separately
supplied value against the charge, so any interpretation as an
operator-norm lower bound needs additional hypotheses and evidence.
Records with different rationales may describe the same endpoint pair;
the charge counts records, not necessarily distinct pairs.

\emph{Computational verification:}
\nolinkurl{code/symbol_game_downmap.py} and
\nolinkurl{code/tests/test_ddc_and_downmap_mutations.py::test_symbol_game_downmap_distinction_collapse_and_loss_ledger}.

\subsubsection{4.7a Flatness does not protect
information}\label{a-flatness-does-not-protect-information}

\begin{quote}
\textbf{{[}PROVEN, EXAMPLE{]} Proposition 4.7a --- Positive Normalized
Flat Loss.} Let every fiber of the three-vertex chain be
\(\mathbb R^3\), and assign every strict comparable pair the map \[
P=\begin{pmatrix}1/2&1/2&0\\1/2&1/2&0\\0&0&1\end{pmatrix}.
\] Then the restriction system is flat, positive,
probability-normalization preserving, and contractive in the induced
\(\ell^1\) norm. Nevertheless, it identifies \(x=e_1\) and \(y=e_2\),
although the protected question \(q(v)=v_1\) has \(q(x)-q(y)=1\).

\emph{Proof.} Direct multiplication gives \(P^2=P\), hence the unique
triple has \(K=P^2-P=0\). Entries are nonnegative and every column sums
to one, so \(P\) preserves the probability simplex and
\(\|P\|_{1\to1}=1\). But \(Px=Py=(1/2,1/2,0)^T\), while \(q(x)=1\) and
\(q(y)=0\). Counting this collapsed protected record gives the declared
charge \(2\), independently of the zero operator norm. \(\square\)
\end{quote}

Consequently, positivity, normalization, contractivity and a positive
lost-record charge cannot imply any strictly positive curvature floor.
The map fails an obligation to preserve \(q\), even though it satisfies
composition flatness. This is the WDP protected-question boundary:
consistency of operations and sufficiency of retained answers are
separate properties. As a supplied quantum example, qubit dephasing also
identifies the two \(\sigma_x\) eigenstates with \(I/2\) and is
idempotent, so using it on each edge again gives \(K=0\). Its projector
Kraus operators certify a completely positive trace-preserving map; this
assumes the qubit carrier and does not derive it from histories.

\subsubsection{4.7b What does give a quantitative curvature
bound}\label{b-what-does-give-a-quantitative-curvature-bound}

\begin{quote}
\textbf{{[}PROVEN, CONDITIONAL{]} Proposition 4.7b --- Measured Path
Separation.} In declared coordinate bases, let \(B=r_{b,a}r_{c,b}\),
\(R=r_{c,a}\), and \(K=B-R\). Suppose an input difference \(d\) and
output linear question \(\ell\) satisfy \(Rd=0\) and
\(|\ell Bd|\ge\eta>0\). Then \[
\|K\|_{1\to1}\ge\frac{\eta}{\|\ell\|_\infty\|d\|_1},\qquad
\sum_{ij}|K_{ij}|\ge\frac{\eta}{\|\ell\|_\infty\|d\|_\infty}.
\] \emph{Proof.} Since \(Rd=0\), \(|\ell Kd|=|\ell Bd|\ge\eta\). The
triangle inequality bounds this by
\(\|\ell\|_\infty\|Kd\|_1\le\|\ell\|_\infty\|K\|_{1\to1}\|d\|_1\).
Alternatively, expanding the finite sum bounds it by
\(\|\ell\|_\infty\|d\|_\infty\sum_{ij}|K_{ij}|\). Positive separation
makes both denominators positive. Division yields both bounds. The
one-coordinate choice \(B=[\eta/(LV)]\), \(R=[0]\), \(d=[V]\),
\(\ell=[L]\), with \(L,V>0\), attains equality. \(\square\)
\end{quote}

This premise compares two actual paths: one destroys the difference,
while the other retains a measured answer. It is stronger than loss
under a single map. Rescaling \(B\) rescales the measured separation and
curvature together. Repeating a record does not multiply either norm. An
additive entrywise bound \(m\eta\) instead requires \(m\) distinct
coordinates with independently verified \(|K_{ij}|\ge\eta\); a charge
convention supplies neither coordinate disjointness nor separation.

\emph{Evidence scope.} \nolinkurl{omega_functional::loss_curvature}
checks the operators, endpoint shapes, direct collapse and positive
separation, and computes both norms from all entries.
\nolinkurl{verify_triangle_separation} constructs \(B\) from the two
supplied restrictions. Independent controls are
\texttt{code/loss\_curvature.py} and
\nolinkurl{code/tests/test_loss_curvature.py}.
\nolinkurl{lean4/FlatLossAndCurvature.lean} formalizes the explicit
countermodel; it does not formalize the general norm inequalities. The
quantum control is a calculation inside its supplied finite carrier.

\subsubsection{4.7c Minimal linear reports and lawful
transport}\label{c-minimal-linear-reports-and-lawful-transport}

Let \(Q\) be a finite rational question matrix. The rows of the nonzero
RREF matrix \(B\) form a basis of its row space. Select its pivot
columns and let \(S\) inject the corresponding coordinate vectors. Then
\[
BS=I,\quad Q=(QS)B,\quad P=I-SB,\quad \operatorname{im}P=\ker B=\ker Q.
\] Indeed the pivot submatrix of \(B\) is \(I\); expressing each row of
\(Q\) in the row basis yields the decoder identity. The converse
inclusion of row spaces comes from row elimination, giving \(Qx=Qy\) iff
\(Bx=By\). Thus \(B\) realizes \(V/\ker Q\), the coarsest quotient
preserving these answers, of dimension \(\operatorname{rank}Q\). Any
other answer-preserving report must distinguish different \(B\)-classes;
a linear such report has rank at least that of \(Q\). All-zero questions
give a zero-dimensional report, not a division error.

For declared \(T:V\to W\) with source and target bases \(B_s,B_t\), a
report map \(D\) with \(DB_s=B_tT\) exists exactly when \[
B_tTP_s=0.
\] Necessity follows by multiplying the replay identity by \(P_s\). For
sufficiency put \(D=B_tTS_s\) and use \(I=S_sB_s+P_s\). It is unique
because any competing replay multiplied by \(S_s\) equals \(D\). For
composable admitted \(T,U\), uniqueness gives \(D(UT)=D(U)D(T)\).
Original target answers replay as \(Q_tT=(Q_tS_t)D(T)B_s\). By induction
this holds for every composition of admitted maps, so the kernel
condition protects the specified future linear experiments as well as
the current report. Equal dimensions or ranks alone are insufficient:
for \(Q=(1,0)\),
\(T=\left(\begin{smallmatrix}1&1\\0&1\end{smallmatrix}\right)\) exposes
the invisible vector \((0,1)\) and is refused.

This is a precise finite linear instance of WDP operational minimality.
It presupposes the state, scalar and question families. It does not
derive positivity, quantum composition, costs, continuation
availability, physical substance, or all native Genesis histories. A
protected probability experiment may be represented by linear questions
only after its probability carrier and admitted operations are
independently specified. The reusable owner is
\nolinkurl{omega_functional::protected_linear_quotient}; Python uses a
different row basis and obtains the same equivalence classes. The
split-descent Lean kernels assume a checked section and kernel premise;
RREF construction and minimality are written proofs and executed exact
algebra.

\subsubsection{4.7d The actual composed-path error
bound}\label{d-the-actual-composed-path-error-bound}

For endpoint-compatible maps \(A_i,B_i:V_i\to V_{i+1}\), acting
chronologically on columns, write \(A_{>i}=A_{m-1}\cdots A_{i+1}\) and
\(B_{<i}=B_{i-1}\cdots B_0\), with identity empty products. Then \[
A_{m-1}\cdots A_0-B_{m-1}\cdots B_0
=\sum_{i=0}^{m-1}A_{>i}(A_i-B_i)B_{<i}.
\] Replace the rightmost differing factor one at a time, or induct on
\(m\); neighboring intermediate products cancel. For compatible induced
norms, \[
\|A_{m-1}\cdots A_0-B_{m-1}\cdots B_0\|
\le\sum_i\|A_{>i}\|\,\|A_i-B_i\|\,\|B_{<i}\|.
\] This follows from the triangle inequality and submultiplicativity and
bounds the final answer discrepancy by the displayed quantity times
\(\|v\|\). The Rust owner \nolinkurl{omega_functional::path_error}
retains both actual paths, computes each term, replays the residual, and
measures the induced \(\ell^1\) norms of the prefix and suffix products.
A sum of local errors is valid without weights only under independently
established contractivity. The scalar pair \((A_0,B_0)=(1,0)\) followed
by \((100,100)\) has local total error 1 and endpoint error 100. For all
\(A_i=2,B_i=1\), the weighted bound is exactly \(2^m-1\), so its
coefficients cannot be uniformly decreased.
\texttt{TransportIdentities.lean} formalizes the arbitrary finite
ring-word expansion; the norm theorem is written, with exact native and
Python controls. No numerical error certificate grants operational
admission or authenticates a declared price. Full native DDC/history
transport remains a separate goal.

\subsubsection{4.8 Simplicial Face-Poset Operator Cocycle (Theorem
4.8)}\label{simplicial-face-poset-operator-cocycle-theorem-4.8}

\begin{quote}
\textbf{{[}PROVEN, CONDITIONAL{]} Theorem 4.8 --- Face-Poset Operator
Identity.} Let \(\mathcal P(T)\) be the finite face-inclusion poset of a
simplicial complex, equipped with declared linear restrictions on
comparable faces. Its composition curvature satisfies \(\delta_2K=0\) on
every four-face chain.

\emph{Proof.} Apply Theorem 3.6 to this finite poset. No metric or
rotation-group assumption is needed. \(\square\)
\end{quote}

The module's Regge terminology is a model interpretation. Its rational
matrix fixtures need not lie in \(\mathrm{SO}(n)\); they contain no
metric edge lengths or dihedral angles. A separately specified ordered
loop product \(W\) has defect \(W-I\), which is not automatically the
comparable-triple defect \(r_{b,a}r_{c,b}-r_{c,a}\). Deriving a Regge
deficit-angle formula or a geometric Bianchi identity requires a
geometric connection and a proved identification; neither follows from
these fixtures.

\emph{Computational verification:}
\nolinkurl{code/tests/test_regge_face_poset.py::test_regge_bianchi_identity_across_dimensions}
verifies this operator identity on the declared face-poset examples.

\subsubsection{4.9 Bell Contextuality and Global Probability Extension
(Theorem
4.9)}\label{bell-contextuality-and-global-probability-extension-theorem-4.9}

In the Abramsky--Brandenburger framework, noncontextuality concerns a
global probability model whose marginals reproduce the observed
contexts. It is distinct from functoriality of marginalization maps.

\begin{quote}
\textbf{{[}PROVEN, CONDITIONAL{]} Theorem 4.9 --- Finite Gluing
Criterion.} Let \(e\) stack the four normalized Bell--CHSH context
distributions, let \(M_{\mathrm{glue}}\) map a distribution on the 16
deterministic assignments to those marginals, and let \(\Delta_{16}\)
denote the nonnegative normalized simplex. Define \[
\Gamma_{\mathrm{context}}(e)=\min_{p\in\Delta_{16}}\|e-M_{\mathrm{glue}}p\|_1.
\] Then \(\Gamma_{\mathrm{context}}(e)=0\) if and only if such a global
probability extension exists. The marginal restriction system itself is
flat, including when \(e\) is a contextual PR box.

\emph{Proof.} The simplex is compact, so the continuous objective
attains its minimum. Its value is zero exactly when all marginal
coordinates agree. Composition of marginalization equals marginalization
to the smaller context, independently of \(e\), so the operator
curvature is zero. \(\square\)
\end{quote}

The exact finite LP gives \(\Gamma_{\mathrm{context}}=2\) and CHSH value
\(S=4\) for the PR box, with zero primal-dual gap. Deterministic local
models have distance zero. The stacked four-context \(L_1\)
normalization is essential: this quantity is not \(\|K\|_1\) and is not
the distance \(D_1\) of §4.3. Tests of the PR/white-noise family at
rational parameters do not constitute a certificate for an arbitrary
irrational quantum model. The algebraic number in §4.10 is enclosed
separately.

\emph{Computational verification:}
\nolinkurl{code/tests/test_quantum_sheaf_contextuality.py::test_pr_box_maximal_contextuality}
and
\nolinkurl{code/tests/test_review_boundaries.py::test_flat_bell_restrictions_do_not_imply_global_probability_extension}.

\subsubsection{4.10 Exact Enclosure of an Algebraic Number (Example
4.10)}\label{exact-enclosure-of-an-algebraic-number-example-4.10}

\begin{quote}
\textbf{{[}PROVEN, EXAMPLE{]} Statement 4.10 --- Rational Enclosure.}
For the algebraic number \(q=2(\sqrt2-1)\), \[
\frac{828}{1000}<q<\frac{829}{1000},\qquad U-L=\frac1{1000}.
\] \emph{Proof.} Both \(1+L/2\) and \(1+U/2\) are positive, and exact
rational squaring gives \[
(1+L/2)^2=\frac{499849}{250000}<2
<\frac{8003241}{4000000}=(1+U/2)^2.
\] Monotonicity of squaring on positive reals yields the enclosure.
\(\square\)
\end{quote}

This establishes an enclosure for this algebraic expression. It does not
identify it with a computed operator curvature, execute an
irrational-input Bell LP, or provide an algorithm for arbitrary
transcendental inputs without an effective representation.

\emph{Computational verification:}
\nolinkurl{code/tests/test_review_boundaries.py::test_algebraic_contextuality_rational_enclosure}
includes a shifted-endpoint negative control.

\subsubsection{4.10a Metric-bound dihedral and deficit
enclosures}\label{a-metric-bound-dihedral-and-deficit-enclosures}

Fix a triangle hinge and a distinct cyclic link of
\(m\in\{3,\ldots,16\}\) vertices. Consecutive link vertices and the
hinge specify the incident 4-simplices. Supply exactly one rational
squared length for every edge used by these simplices; shared faces
therefore inherit the same metric. Each simplex must have a
positive-definite Euclidean Gram matrix. These are finite local metric
and gluing data, without a global spacetime or a comparison to an
arbitrary supplied holonomy word.

Using one hinge vertex as origin, form the \(4\times4\) rational Gram
matrix \[
G_{ij}=(s_{0i}+s_{0j}-s_{ij})/2.
\] Positive leading principal minors certify positive definiteness.
Separate the two tangent hinge directions from the two opposite vertices
and take the Schur complement. Write its entries as \(a,b,c\), with
\(a,b>0\) and \(d=ab-c^2>0\). The interior dihedral angle is \[
\theta=\begin{cases}
\arctan(\sqrt d/c),&c>0,\\
\pi/2,&c=0,\\
\pi-\arctan(\sqrt d/(-c)),&c<0.
\end{cases}
\] This follows by orthogonally removing the hinge-tangent components:
the remaining normal-plane vectors have dot product \(c\), squared norms
\(a,b\), and unsigned parallelogram area \(\sqrt d\). The branch is the
interior angle in \((0,\pi)\).

\begin{quote}
\textbf{{[}PROVEN, CONDITIONAL{]} Proposition 4.10a --- Certified Finite
Geometric Enclosure.} For these effective inputs, the following rational
procedures enclose each specified dihedral angle and the deficit
\(\delta=2\pi-\sum_i\theta_i\). Each accepted result has width at most
its positive rational request \(\varepsilon\). Unsupported inputs and
exhausted arithmetic budgets receive a refusal, not a sign verdict.

\emph{Proof.} For \(q\ge0\), choose \(b\) with \(2^{-b}\le\varepsilon\)
and let \(k=\lfloor\sqrt{\lfloor q2^{2b}\rfloor}\rfloor\). Then
\(k^2\le q2^{2b}<(k+1)^2\). Thus \([k/2^b,(k+1)/2^b]\) encloses
\(\sqrt q\) with the required width; rational square roots may be
returned exactly.

For \(0\le x\le1/2\), integrate the finite geometric identity for
\(1/(1+t^2)\) to obtain \[
\arctan x=\sum_{k=0}^{N-1}\frac{(-1)^kx^{2k+1}}{2k+1}
+(-1)^N\int_0^x\frac{t^{2N}}{1+t^2}\,dt.
\] The integral lies between zero and \(x^{2N+1}/(2N+1)\le4^{-N}\). The
sum and its signed next-term endpoint therefore enclose the answer.
Oddness handles negative inputs. For \(1/2<x\le1\), use
\[\arctan x=\arctan(1/2)+\arctan\frac{x-1/2}{1+x/2},\] whose second
argument lies in \([0,1/3]\). For \(x>1\), use
\(\arctan x=\pi/2-\arctan(1/x)\).

Machin's identity \(\pi=16\arctan(1/5)-4\arctan(1/239)\) supplies a
rational enclosure of \(\pi\). To check the identity, successive
tangent-doubling gives \(5/12\) and \(120/119\); subtracting the angle
with tangent \(1/239\) yields tangent one. The angle
\(4\arctan(1/5)-\arctan(1/239)\) lies in \((0,\pi/2)\), so it equals
\(\pi/4\). The elementary bounds \(x-x^3/3\le\arctan x\le x\) establish
the required branch. Budget each arctangent width at \(\varepsilon/20\)
to obtain a \(\pi\) interval of width at most \(\varepsilon\).

For \(c\ne0\), set \(t=\varepsilon/4\) and enclose \(\sqrt d\) with
width at most \(|c|t\). Evaluate arctangent at both rational endpoints
divided by \(|c|\), each with width at most \(t\). Monotonicity and
\(|\arctan u-\arctan v|\le|u-v|\) bound the combined width by \(3t\).
The obtuse branch adds a \(\pi\) interval of width at most \(t\). The
right-angle branch uses \(\pi/2\). Finally, for a star, request each
angle and \(\pi\) to width \(\varepsilon/(m+2)\). Rational interval
subtraction gives total width at most \(\varepsilon\). \(\square\)
\end{quote}

The implementation profile checks positive denominators and at most
1,024 numerator/denominator bits, permits 1,024 series terms and 2,048
dyadic bits, and refuses an intermediate argument exceeding its profile.
For accepted scalar input requests, \(\varepsilon\ge2^{-1024}\); all
series sub-budgets need fewer than 1,024 terms, and dyadic refinement
needs at most 1,024 bits.

A simple geometry class with guaranteed resource termination has integer
squared lengths between \(1\) and \(2^{16}\), positive Gram matrices,
and \(0<\varepsilon\le1\) with numerator and denominator at most 64
bits. The hinge tangent determinant numerator is a positive integer at
most \(2^{34}\). Schur entries have a common denominator at most
\(2^{35}\), absolute value at most \(2^{16}\), and nonzero
\(|c|\ge2^{-35}\). Including the star error budget needs fewer than 110
dyadic bits. The discriminant and arctangent endpoint numerators and
denominators consequently stay below 1,024 bits. More general rational
metric inputs are evaluated when their intermediates fit the declared
profile; validity of a geometry alone is not a promise to bypass
resource refusals.

\emph{Independent reference.} \nolinkurl{code/certified_geometry.py}
uses bisection for square roots and a different positive series: \[
\arctan x=\frac{x}{1+x^2}\sum_{k\ge0}c_kz^k,\quad
z=\frac{x^2}{1+x^2},\quad c_0=1,\quad c_k=\frac{2k}{2k+1}c_{k-1}.
\] For \(0\le x\le1\), expand the integral
\(x\int_0^1(1+x^2t^2)^{-1}\,dt\) in powers of \(z(1-t^2)\). Integration
by parts gives these coefficients, with \(0<c_k\le1\); the tail after
\(N\) terms is at most \(xz^N\). The reference obtains \(\pi\) from
\(4\arctan1\), independently of Machin's identity, and uses SymPy
determinants and inversion rather than the native Gaussian routine.

\subsubsection{4.10b A sign verdict needs separation
evidence}\label{b-a-sign-verdict-needs-separation-evidence}

An interval \([L,U]\) with \(L>0\) certifies \(\delta\ge L>0\); one with
\(U<0\) certifies \(\delta\le U<0\). These endpoints are verified
separation evidence computed from the metric. An interval containing
zero returns \texttt{ContainsZero} and terminates without a nonvanishing
claim. Refinement may be repeated, but no bound ensures zero exclusion
on flat inputs. If a separate proof establishes \(|\delta|\ge\sigma>0\),
any containing interval of width less than \(\sigma\) must exclude zero.
An unverified caller assertion of separation is insufficient.

Regular 4-simplices give \(\cos\theta=1/4\): three- and four-sector
stars have positive deficits, while five- and six-sector stars have
negative deficits. Four orthogonal sectors have \(\theta=\pi/2\) and
zero deficit, retained by the interval sign classifier. These are
controls within the declared input class. They do not identify geometric
\(\delta\) with comparable-triple operator \(K\), nor prove an
unrestricted geometry or quantum theorem.

\emph{Evidence scope.} Reusable scalar enclosures belong to
\nolinkurl{omega_number_theory::certified_angles}, reusing
\texttt{omega\_kernel::interval} and the existing integer square root.
Metric extraction and the source-bound closed-star report belong to
\nolinkurl{omega_general_relativity::certified_dihedral}, reusing
\nolinkurl{linear_algebra::regge_exact}. Focused native and independent
Python results are recorded in
\href{docs/loss_incidence_angles_2026_10_14.md}{the current derivation
review}. Containment and stopping are written mathematical proofs; no
Lean proof of the transcendental estimates is claimed.

\subsubsection{4.10c Metric hinge transport and its exact
observation}\label{c-metric-hinge-transport-and-its-exact-observation}

A Euclidean metric hinge star from §4.10a supplies its cyclic sectors
and positive normal Gram data \((a_i,b_i,c_i)\). In each hinge-adapted
oriented orthonormal frame the consecutive normal rays have interior
angle \(\theta_i\), with \[
\cos\theta_i=\frac{c_i}{\sqrt{a_ib_i}},\qquad
\sin\theta_i=\frac{\sqrt{a_ib_i-c_i^2}}{\sqrt{a_ib_i}}>0.
\] Declare face gluing that fixes the identified hinge tangent plane and
changes normal coordinates by \(R(-\theta_i)\). This is the oriented
unfolding connection of the piecewise Euclidean star: the metric
determines each angle, while the common hinge identification, cyclic
orientation and adapted frames are additional declared data. The links
are \(H_i=\operatorname{diag}(I_2,R(-\theta_i))\). Their chronological
loop is \[
W=H_{m-1}\cdots H_0
=\operatorname{diag}(I_2,R(-\sum_i\theta_i))
=\operatorname{diag}(I_2,R(\delta)),\quad
\delta=2\pi-\sum_i\theta_i.
\] Rotations add their angles; \(2\pi\) periodicity yields the last
equality. Reversing the cyclic path reverses order and inverts each
link, giving \(W^{-1}\). At local frame changes \(g_i\in SO(4)\) the
links become \(H_i'=g_{i+1}H_ig_i^{-1}\), with \(g_m=g_0\). Telescoping
gives the loop at a specified common basepoint \(b\) as
\(W_b'=g_bW_bg_b^{-1}\). Products at different basepoints must be
compared by their connecting transport, not silently multiplied as if
their frames agreed.

The exact Rust class accepts precisely stars for which these normal
cosines and sines are rational; it tests rational square roots of their
squares, including the sign of \(c_i\), and refuses other algebraic
rotations. Its metric angle intervals remain valid beyond this class. A
declared triangle observation sets its early map to \(H_0\), late map to
\(H_{m-1}\cdots H_1\) and direct map to \(I\) after closed
endpoint-frame identification. Only for this explicit observation does
\(K=W-I\). It is not an identification for arbitrary posets.

Two countermodels fix the interpretation. Three orthogonal sectors have
\(\delta=\pi/2\) and \(W\ne I\); an explicit four-vertex operator
extension still satisfies the nonempty quadruple identity \(DK=0\). Thus
additive Bianchi alone does not force this ordered metric loop to close.
Eight orthogonal sectors have \(\delta=-2\pi\) yet \(W=I\). Thus even
this metric correspondence loses the unwrapped deficit unless
winding/branch information is retained. No Lorentzian, global spacetime,
arbitrary algebraic rotation or quantum physical reconstruction is
claimed. Metric, oriented connection, basepoint, observation and
operational availability remain separate premises. The owner is
\nolinkurl{omega_general_relativity::metric_hinge_transport}; its forged
frame and metric controls independently revalidate public matrix inputs.

\subsubsection{4.11 Exact Non-Abelian Gauge Holonomies on 4D Simplicial
Complexes (Theorem
4.11)}\label{exact-non-abelian-gauge-holonomies-on-4d-simplicial-complexes-theorem-4.11}

The following finite matrix construction studies non-abelian ordered
holonomies. It supplies exact algebraic fixtures; it makes no claim of
novelty over all prior discrete gauge formulations or of a physical
spacetime reconstruction.

We implement exact discrete non-abelian gauge group representations over
the cyclotomic field \(\mathbb{Q}(\zeta_8)\), where
\(\zeta_8 = e^{2\pi i / 8} = \frac{1+i}{\sqrt{2}}\). Because
\(\sqrt{2} = \zeta_8 - \zeta_8^3\) and \(i = \zeta_8^2\), the field
\(\mathbb{Q}(\zeta_8) \cong \mathbb{Q}[x] / \langle x^4 + 1 \rangle\)
furnishes an exact 4-dimensional rational vector space containing all
roots needed for the quaternion group \(Q_8 \subset \mathrm{SU}(2)\),
the binary tetrahedral group \(2T \subset \mathrm{SU}(2)\), the standard
single-qubit Clifford \(+ T\) group \(\langle H, T \rangle\), and exact
4D spacetime rotations in \(\mathrm{SO}(4)\) via the canonical
double-cover epimorphism
\(\mathrm{Spin}(4) \cong \mathrm{SU}(2) \times \mathrm{SU}(2) \to \mathrm{SO}(4)\)
acting by quaternion conjugation \(q \mapsto U_L q U_R^\dagger\).

\begin{quote}
\textbf{{[}PROVEN, CONDITIONAL{]} Theorem 4.11 --- Exact Non-Abelian 4D
Simplicial Gauge Holonomy.} Let \(\Delta_4\) be a 4-dimensional
simplicial complex and let \(h = \sigma^2\) be a 2-dimensional
triangular hinge shared by a cyclic star of \(m\) 4-simplices
\((\sigma^4_0, \sigma^4_1, \dots, \sigma^4_{m-1})\). Assigning exact
unitary link holonomies \(U_e \in \mathcal{U}_n(\mathbb{Q}(\zeta_8))\)
to directed dual edges connecting adjacent 4-simplices around \(h\): 1.
\textbf{Path-Ordered Wilson Holonomy:} The hinge holonomy
\(W(h) = \mathcal{P} \prod_{i=0}^{m-1} U_{\sigma^4_i, \sigma^4_{i+1}}\)
is strictly unitary (\(W(h)^\dagger W(h) = I\)) and evaluated in exact
rational arithmetic over \(\mathbb{Q}(\zeta_8)\). 2. \textbf{Curvature
Defect Operator:} The simplicial non-abelian curvature defect
\(K(h) = W(h) - I\) vanishes (\(K(h) = 0\)) if and only if the dual
connection around the hinge is gauge-flat. 3. \textbf{Gauge Covariance
\& Trace Invariance:} Under local discrete gauge transformations
\(U_{ij} \mapsto \Omega_i U_{ij} \Omega_j^\dagger\) with
\(\Omega_i \in \mathcal{U}_n(\mathbb{Q}(\zeta_8))\), the Wilson holonomy
transforms by conjugation
\(W(h) \mapsto \Omega_0 W(h) \Omega_0^\dagger\), ensuring that the
normalized Wilson loop trace \(\frac{1}{n}\operatorname{Tr}(W(h))\) is
strictly gauge-invariant in \(\mathbb{Q}(\zeta_8)\). 4.
\textbf{Quaternion \(\mathrm{SO}(4)\) Action:} For any pair of
\(\mathrm{SU}(2)\) holonomies \((U_L, U_R)\) over
\(\mathbb{Q}(\zeta_8)\), the induced map on
\(\mathbb{R}^4 \cong \mathbb{H}\) produces an exact orthogonal matrix
\(R \in \mathrm{SO}(4) \cap \mathcal{M}_4(\mathbb{Q}(\sqrt{2}))\)
satisfying \(R^T R = I_4\) and \(\det(R) = +1\) with zero rounding
error.

\emph{Proof.} Exact polynomial arithmetic in
\(\mathbb{Q}[x]/\langle x^4+1 \rangle\) satisfies the ring and field
axioms without rounding error. Dual basis inversion and unitarity are
certified by exact conjugate transposition. Gauge covariance follows
from the cyclic cancellation \(\Omega_i^\dagger \Omega_i = I\) in the
path-ordered product. The quaternion action gives the standard
double-cover homomorphism to \(\mathrm{SO}(4)\) in the \(\{1,i,j,k\}\)
basis; \((-I,-I)\) acts trivially, so this is not an injective
embedding. \(\square\)
\end{quote}

\emph{Computational verification:} Implemented in
\texttt{code/cyclotomic\_gauge.py} and tested in
\nolinkurl{code/tests/test_cyclotomic_gauge_4d.py}.

\begin{quote}
\textbf{{[}PROVEN, CONDITIONAL{]} Corollary 4.11a --- Discrete
Non-Abelian Bianchi Identity on 4-Simplices.} On any 4-simplex
\(\sigma^4 = (v_0, v_1, v_2, v_3, v_4)\) with unitary link variables
\(U_{uv} \in \mathcal{U}_n(\mathbb{Q}(\zeta_8))\) assigned to its 10
edges: 1. \textbf{Ordered Tetrahedral Product Identity:} For each 3-cell
(tetrahedron) \(\tau = (v_0, v_1, v_2, v_3) \subset \partial \sigma^4\)
with basepoint \(v_0\), the ordered product of the three meeting face
holonomies \(W_1 = W(v_0, v_3, v_1)\), \(W_2 = W(v_0, v_2, v_3)\),
\(W_3 = W(v_0, v_1, v_2)\) and the basepoint-transported opposite face
holonomy
\(W_{\text{opp}, 0} = U_{v_1, v_0} W(v_1, v_2, v_3) U_{v_0, v_1}\)
closes to identity identically: \[
   \mathcal{B}(\tau, v_0) = W_1 \cdot W_2 \cdot W_3 \cdot W_{\text{opp}, 0}^\dagger = I_n.
   \] 2. \textbf{Boundary 4-Simplex Bianchi Closure:} Transporting all 5
tetrahedral Bianchi operators to a common basepoint \(v_0\) via link
variables, the boundary product over \(\partial \sigma^4\) satisfies: \[
   \mathcal{B}(\partial \sigma^4) = \prod_{p=0}^4 \mathcal{B}_0(\tau_p) = I_n,
   \] with zero curvature defect in exact cyclotomic arithmetic over
\(\mathbb{Q}(\zeta_8)\) for arbitrary non-abelian connections (including
\(Q_8\), \(2T\), and \(\mathrm{SO}(4)\)). 3. \textbf{Fixture Mutation
Sensitivity:} Selected noncommuting fixtures produce a nonzero residual
after an orientation, order, or link-consistency mutation. This is a
finite test result, not an unrestricted claim: an identity face loop or
commuting face holonomies can make such a mutation invisible.

\emph{Proof.} With the declared orientation and multiplication order,
cancellation of inverse links gives \(W_1W_2W_3=W_{\mathrm{opp},0}\), so
multiplication by \(W_{\mathrm{opp},0}^\dagger\) gives identity.
Conjugating each identity to the common basepoint and multiplying the
five identities gives the boundary statement. This latter product is a
consequence of the tetrahedral identities, not an independent
higher-coherence theorem. The additive residual
\(F_1+F_2+F_3-F_{\mathrm{opp},0}\) need not vanish; the exact identity
is multiplicative. \(\square\)
\end{quote}

\emph{Computational verification:} Implemented in
\texttt{code/cyclotomic\_gauge.py} (\texttt{FourSimplexGaugeSystem}) and
tested in \nolinkurl{code/tests/test_cyclotomic_gauge_4d.py} (24 passing
unit and adversarial mutation tests).

\begin{center}\rule{0.5\linewidth}{0.5pt}\end{center}

\subsection{5. Discussion and Structural
Comparison}\label{discussion-and-structural-comparison}

\begin{longtable}[]{@{}
  >{\raggedright\arraybackslash}p{(\linewidth - 6\tabcolsep) * \real{0.2500}}
  >{\raggedright\arraybackslash}p{(\linewidth - 6\tabcolsep) * \real{0.2500}}
  >{\raggedright\arraybackslash}p{(\linewidth - 6\tabcolsep) * \real{0.2500}}
  >{\raggedright\arraybackslash}p{(\linewidth - 6\tabcolsep) * \real{0.2500}}@{}}
\toprule\noalign{}
\begin{minipage}[b]{\linewidth}\raggedright
Property
\end{minipage} & \begin{minipage}[b]{\linewidth}\raggedright
Path A
\end{minipage} & \begin{minipage}[b]{\linewidth}\raggedright
Path B
\end{minipage} & \begin{minipage}[b]{\linewidth}\raggedright
Finite extensions (Path C)
\end{minipage} \\
\midrule\noalign{}
\endhead
\bottomrule\noalign{}
\endlastfoot
Input & Cover maps with path-independent saturated chains & Declared
comparable-pair maps & Declared maps, bases, and model-specific data \\
Output & Unique strict presheaf & Composition defect \(K\) & Conditional
additive normal form; finite LP certificates; routing traces \\
Differential & \(d_1d_0=0\) & \(d_1d_0s=K(s)\) & \(\delta_2K=0\) and the
stated curvature residual for \(\delta_2\delta_1\) \\
Categorical status & Strict 1-categorical interface & No generic \(H^2\)
quotient & All-degree strict-chain curved algebra; full lax 2-presheaf
remains open \\
Diagnostics & Flatness under the cover hypothesis & Local compatibility
can coexist with nonzero \(K\) & Distance-to-image, probability-gluing
distance, and ledger charges have separate definitions \\
\end{longtable}

\subsection{6. Scope and Open
Obligations}\label{scope-and-open-obligations}

The claims concern finite posets, specified linear maps, and the stated
finite matrix and LP models. They do not establish physical dynamics,
infinite-poset limits, observer-independent norms, or higher cohomology.
A dual LP certificate does not repair the missing inclusion
\(\operatorname{im}\delta_1\subseteq\ker\delta_2\) and cannot promote a
non-flat system to a generic \(H^2\) class.

The five connected derivations now supply faithful heterogeneous-fiber
Lean transport, a specified observed 2-category and normal lax functor,
a metric hinge-holonomy correspondence for the exact rational rotation
class, a minimal protected linear report with unique lawful transport,
and sharp weighted path-error bounds. Their current scopes and
executions are recorded in
\href{docs/five_connected_derivations_2026_10_14.md}{the five-derivation
supplement}. Graded Jacobi remains written and native rather than
formally replayed; the typed 2-category is not wholly Lean-formalized.
General geometric connections and algebraic/Lorentzian extensions remain
open. The full native Genesis/DDC correspondence is blocked on
preservation/reflection, continuations, prices and availability. Section
4.4 classifies only the declared vertexwise invertible basis-change
action. Operator flatness, loss of information and Bell positive gluing
remain distinct; none forces a quantum carrier or identifies physical
substance.

\subsection{Appendix A: Mechanism
Inventory}\label{appendix-a-mechanism-inventory}

The CODA terms below name programme roles. Their mathematical content is
limited to the stated definitions and hypotheses.

\begin{longtable}[]{@{}
  >{\raggedright\arraybackslash}p{(\linewidth - 2\tabcolsep) * \real{0.5000}}
  >{\raggedright\arraybackslash}p{(\linewidth - 2\tabcolsep) * \real{0.5000}}@{}}
\toprule\noalign{}
\begin{minipage}[b]{\linewidth}\raggedright
Slot
\end{minipage} & \begin{minipage}[b]{\linewidth}\raggedright
Value
\end{minipage} \\
\midrule\noalign{}
\endhead
\bottomrule\noalign{}
\endlastfoot
Primitive differences & Cover-edge section differences; payload
distinctions; declared protected distinction records. \\
Carriers / structure & Finite posets, linear fibers, comparable-pair
maps, finite face posets, and Bell measurement contexts. \\
Cochain sequence & Section-valued \(C^0,C^1,C^2,C^3\) and the specified
operator-cochain spaces. Non-flat operators need not form a complex. \\
Coboundary maps & The declared \(d_0,d_1,d_2\) and
\(\delta_1,\delta_2\), with the identities of §3 and §4.1. \\
Zero relation & Operator flatness \(K=0\); separately, equal payloads
and zero diamond path difference. \\
Licensed repairs & Unique strict completion under path independence;
conditional additive normal form; finite rational distance-to-image
optimization. \\
Quotient / Q1 & Payload quotient only when its classifier descends. No
generic non-flat cohomology quotient is constructed. \\
Residuals & Composition curvature; section-valued degree-3 boundary;
operator differential residual; separately declared ledger charge
\(2m\). \\
Witnesses & Exact Python fixtures and primal-dual pairs; selected
current Lean elaboration and declaration-specific axiom output;
historical manifests remain separate. \\
Observers / profiles & Section and payload evaluators, declared
coordinate norms, probability-gluing LP, and finite routing profile. \\
Anomalons & Residual native-history, general geometric, graded-Jacobi
formalization and physical reconstruction obligations; exact
counterexamples prevent stronger promotion. \\
Completion law & Strict chain composition under the cover hypothesis;
\(d_1d_0s=K(s)\); \(\delta_2K=0\). \\
Terminal / route & Choose Path A under its hypothesis or record Path B
defects; then apply each finite extension under its own hypotheses.
Evidence is claim-specific. \\
Authority status & None. These mathematical and model statements carry
no operational or institutional authority. \\
\end{longtable}

\begin{center}\rule{0.5\linewidth}{0.5pt}\end{center}

\subsection{Appendix B: Claim-Type
Ledger}\label{appendix-b-claim-type-ledger}

The IDs match both companion ledgers. Prose statements, finite tests and
the selected current Lean replay have distinct scopes. Historical source
presence does not certify a current whole-package build.

\begin{longtable}[]{@{}
  >{\raggedright\arraybackslash}p{(\linewidth - 8\tabcolsep) * \real{0.2000}}
  >{\raggedright\arraybackslash}p{(\linewidth - 8\tabcolsep) * \real{0.2000}}
  >{\raggedright\arraybackslash}p{(\linewidth - 8\tabcolsep) * \real{0.2000}}
  >{\raggedright\arraybackslash}p{(\linewidth - 8\tabcolsep) * \real{0.2000}}
  >{\raggedright\arraybackslash}p{(\linewidth - 8\tabcolsep) * \real{0.2000}}@{}}
\toprule\noalign{}
\begin{minipage}[b]{\linewidth}\raggedright
\#
\end{minipage} & \begin{minipage}[b]{\linewidth}\raggedright
ID
\end{minipage} & \begin{minipage}[b]{\linewidth}\raggedright
Claim
\end{minipage} & \begin{minipage}[b]{\linewidth}\raggedright
Type
\end{minipage} & \begin{minipage}[b]{\linewidth}\raggedright
Evidence / scope
\end{minipage} \\
\midrule\noalign{}
\endhead
\bottomrule\noalign{}
\endlastfoot
1 & PO001-C01 & Definition 2.1: Functorial Banach Presheaf &
\texttt{{[}DEFINITION{]}} & §2.1; definitional; not applicable. \\
2 & PO001-C02 & Proposition 2.1a: Strict Completion from Covering Edges
& \texttt{{[}PROVEN,\ CONDITIONAL{]}} & §2.1; prose proof; lean source
present not replayed. \\
3 & PO001-C03 & Lemma 2.2: Diamond Path-Independence &
\texttt{{[}PROVEN,\ CONDITIONAL{]}} & §2.2; prose proof; lean source
present not replayed. \\
4 & PO001-C04 & Proposition 2.3: Categorical Interface Unlock &
\texttt{{[}STANDARD,\ CONDITIONAL{]}} & §2.3; standard theory; not
applicable. \\
5 & PO001-C05 & Definition 3.1: Composition Curvature Operator &
\texttt{{[}DEFINITION{]}} & §3.1; definitional; lean source present not
replayed. \\
6 & PO001-C06 & Theorem 3.2: Coboundary-Curvature Identity d1 d0 = K &
\texttt{{[}PROVEN,\ CONDITIONAL{]}} & §3.2; prose proof; lean source
present not replayed. \\
7 & PO001-C07 & Corollary 3.3: Flat Locus Equivalence &
\texttt{{[}PROVEN,\ CONDITIONAL{]}} & §3.2; prose proof; lean source
present not replayed. \\
8 & PO001-C08 & Scalar Diamond Curvature Witness &
\texttt{{[}PROVEN,\ EXAMPLE{]}} & §2.2; exact calculation; lean source
present not replayed. \\
9 & PO001-C09 & Theorem 3.4: Finite d2 K Boundary Formula &
\texttt{{[}PROVEN,\ CONDITIONAL{]}} & §3.3; prose proof; lean source
present not replayed. \\
10 & PO001-C10 & Theorem 3.5: Payload-Curvature LGGR Boundary &
\texttt{{[}PROVEN,\ CONDITIONAL{]}} & §3.4; prose proof; lean source
present not replayed. \\
11 & PO001-C11 & Theorem 3.6: Unconditional Operator Cocycle Identity &
\texttt{{[}PROVEN,\ CONDITIONAL{]}} & §3.5; prose proof; lean source
present not replayed. \\
12 & PO001-C12 & Theorem 3.7: Generic H\^{}2 Quotient Obstruction &
\texttt{{[}PROVEN,\ EXAMPLE{]}} & §3.5; exact calculation; finite python
fixtures. \\
13 & PO001-C13 & Curvature Cochain Interpretation &
\texttt{{[}INTERPRETIVE{]}} & §3.1, §5, §3.5; discursive; not
applicable. \\
14 & PO001-C14 & Scope Fences & \texttt{{[}SCOPE\ FENCE{]}} & §6;
axiomatic negative; not applicable. \\
15 & PO001-C15 & Open H\^{}2 Promotion & \texttt{{[}OPEN{]}} & §3.5, §6;
open; not mechanized. \\
16 & PO001-C16 & Theorem 4.1: Curved Differential Derivation Identity &
\texttt{{[}PROVEN,\ CONDITIONAL{]}} & §4.1, §4.1a--b; written all-degree
proof; faithful varying-Hom poset Lean transport; graded Jacobi
written/native; bounded exact controls. \\
17 & PO001-C17 & Theorem 4.2: Four-Term Operator Comparison &
\texttt{{[}PROVEN,\ CONDITIONAL{]}} & §4.2--a; written observed
2-category, units and pentagon; ring kernels and exact controls. \\
18 & PO001-C18 & Theorem 4.3: Exact Primal-Dual Distance-to-Image
Certification & \texttt{{[}PROVEN,\ CONDITIONAL{]}} & §4.3; prose proof;
finite python fixtures. \\
19 & PO001-C19 & Theorem 4.4: Cover-Preserving Additive Normal Form &
\texttt{{[}PROVEN,\ CONDITIONAL{]}} & §4.4; prose proof; finite python
fixtures. \\
20 & PO001-C20 & Theorem 4.5: Prime Triad Bounded-Curvature Routing Law
& \texttt{{[}PROVEN,\ CONDITIONAL{]}} & §4.5; prose proof; finite python
fixtures. \\
21 & PO001-C21 & Theorem 4.6: DDC Operational Commutator on Poset
Lattices & \texttt{{[}PROVEN,\ CONDITIONAL{]}} & §4.6; prose proof;
finite python fixtures. \\
22 & PO001-C22 & Statement 4.7: Declared Cost-Two Ledger Charge &
\texttt{{[}MODEL\ CONSEQUENCE{]}} & §4.7--d; declared charge, flat-loss
refutation, protected linear quotient and sharp weighted path bounds. \\
23 & PO001-C23 & Theorem 4.8: Simplicial Face-Poset Operator Cocycle &
\texttt{{[}PROVEN,\ CONDITIONAL{]}} & §4.8; prose proof; finite python
fixtures. \\
24 & PO001-C24 & Theorem 4.9: Bell Global Probability Gluing Criterion &
\texttt{{[}PROVEN,\ CONDITIONAL{]}} & §4.9; prose proof; finite python
fixtures. \\
25 & PO001-C25 & Example 4.10: Rational Enclosure of an Algebraic Number
& \texttt{{[}PROVEN,\ EXAMPLE{]}} & §4.10--c; metric angle enclosures,
rational normal transport and branch countermodel. \\
26 & PO001-C26 & Theorem 4.11: Exact Non-Abelian 4D Simplicial Gauge
Holonomy & \texttt{{[}PROVEN,\ CONDITIONAL{]}} & §4.11; prose proof;
finite python fixtures. \\
27 & PO001-C27 & Corollary 4.11a: Ordered Tetrahedral Product Identity &
\texttt{{[}PROVEN,\ CONDITIONAL{]}} & §4.11; prose proof; finite python
fixtures. \\
\end{longtable}

\begin{center}\rule{0.5\linewidth}{0.5pt}\end{center}

\subsection{Data and Code
Availability}\label{data-and-code-availability}

Paper sources and independent finite reference artifacts are maintained
in this package tree under \texttt{CC-BY-4.0} (text) and
\texttt{Apache-2.0} (code). The working October 14 supplement has not
been redeposited. Its reusable Rust owner is maintained separately in
the repository:

\begin{itemize}
\tightlist
\item
  \textbf{Manuscript:} \texttt{covering\_edge\_descent.md},
  \nolinkurl{covering_edge_descent.tex},
  \nolinkurl{covering_edge_descent.pdf}.
\item
  \textbf{Claim Ledgers:} Companion
  \nolinkurl{covering_edge_descent_claim_ledger.md} and
  \nolinkurl{covering_edge_descent_claim_ledger.yaml}.
\item
  \textbf{Change Log:} Companion
  \nolinkurl{covering_edge_descent_change_log.md}.
\item
  \textbf{Lean 4 Sources (\texttt{lean4/}):} 20 theorem/module sources
  plus \texttt{lakefile.lean}. The current focused replay uses
  explicitly migrated stable Lean 4.34.0 and exact Mathlib/dependency
  pins; the concrete incidence, flat-loss, heterogeneous incidence and
  transport identity modules are elaborated. This does not certify every
  older root or the whole manuscript. The September 27 manifest and
  earlier receipts remain historical records; see
  \texttt{lean4/README.md} and the current derivation review for exact
  limits.
\item
  \textbf{Python Computational Engine (\texttt{code/},
  \texttt{code/tests/}):} Finite exact rational and cyclotomic examples,
  counterexamples, LP certificates, and model traces. Current replay
  results and limitations are recorded in
  \nolinkurl{docs/review_2026_10_02.md}; tests do not establish every
  general mathematical claim.
\item
  \textbf{Reusable Rust:} \nolinkurl{omega_homology::restriction_dga},
  with focused tests in that crate; the standalone paper package
  includes the independent Python model and the source/replay
  commitments in \nolinkurl{code/operator_incidence_verification.json}.
\item
  \textbf{Checksum Manifest:} \texttt{SHA256SUMS} covering all package
  files.
\end{itemize}

Reproduction commands from package root:

\begin{Shaded}
\begin{Highlighting}[]
\ExtensionTok{python}\NormalTok{ scripts/verify\_package.py}
\CommentTok{\# Optional reader{-}format rebuild (Pandoc and XeLaTeX required):}
\ExtensionTok{python}\NormalTok{ scripts/build\_paper.py}
\end{Highlighting}
\end{Shaded}

\begin{center}\rule{0.5\linewidth}{0.5pt}\end{center}

\subsection{References}\label{references}

\begin{enumerate}
\def\labelenumi{\arabic{enumi}.}
\tightlist
\item
  Watrous, J. (2018). \emph{The Theory of Quantum Information}.
  Cambridge University Press.
\item
  Mac Lane, S. (1998). \emph{Categories for the Working Mathematician}.
  Springer.
\item
  Carroll, J. H. (2026). Paper 000, \emph{Projection Obstruction for
  Factorized States}. Zenodo. \url{https://zenodo.org/records/21752743}.
\item
  Carroll, J. H. (2026). Paper 001, \emph{From Covering-Edge Descent
  Complexes to Functorial Presheaves and Coherence Defects}. Zenodo.
  \url{https://zenodo.org/records/20710765}.
\item
  Regge, T. (1961). \emph{General relativity without coordinates}. Il
  Nuovo Cimento, 19(3), 558-571.
\item
  Abramsky, S., \& Brandenburger, A. (2011). \emph{The sheaf-theoretic
  structure of non-locality and contextuality}. New Journal of Physics,
  13(11), 113036. \url{https://arxiv.org/abs/1102.0264}.
\item
  Fine, A. (1982). \emph{Hidden variables, joint probability, and the
  Bell inequalities}. Physical Review Letters, 48(5), 291.
\item
  Loday, J.-L. (1998). \emph{Cyclic Homology}. Springer.
\end{enumerate}

\end{document}
